% Luc Destombe --- licence CC BY-NC-SA 4.0
% https://creativecommons.org/licenses/by-nc-sa/4.0/deed.fr
% Fichier autonome : compiler avec pdflatex, deux fois (signets, nombre de pages).
\documentclass[11pt,a4paper]{article}
\makeatletter
\newif\iflycee@unecolonne
\lycee@unecolonnetrue
\newif\iflycee@palatino
\lycee@palatinotrue

\RequirePackage[utf8]{inputenc}
\RequirePackage[T1]{fontenc}
\RequirePackage[french]{babel}
\RequirePackage{etoolbox}

\newcommand{\ShorthandsOff}{\@ifundefined{shorthandoff}{}{\shorthandoff{;:!?}}}
\newcommand{\ShorthandsOn}{\@ifundefined{shorthandon}{}{\shorthandon{;:!?}}}
\ShorthandsOff
\AtBeginDocument{\ShorthandsOn}
\AtBeginEnvironment{tikzpicture}{\ShorthandsOff}

\RequirePackage{geometry}
\RequirePackage{fancyhdr}
\RequirePackage{lastpage}
\RequirePackage{multicol}
\RequirePackage{array}
\RequirePackage{enumitem}
\RequirePackage{xcolor}
\RequirePackage{needspace}

\RequirePackage{amsmath,amssymb}
\RequirePackage{mathpazo}
\DeclareMathSizes{10.95}{10.95}{8}{5.5}
\begingroup\catcode`\_=\active \catcode`\^=\active
\gdef_#1{\sb{\scriptscriptstyle #1}}
\gdef^#1{\sp{\scriptscriptstyle\lycee@moinsepais #1}}
\endgroup
\newcommand\lycee@moins{\mathbin{%
  \setbox\z@\hbox{$\m@th\scriptscriptstyle\lycee@moinsorig$}%
  \dimen@=.55\wd\z@ \advance\dimen@-.5pt
  \hbox to.75\wd\z@{\hss$\m@th\scriptscriptstyle\vcenter{\hbox{%
    \tikz\draw[line width=.5pt,line cap=round](0,0)--(\the\dimen@,0);}}$\hss}}}
\begingroup\lccode`\~=`\- \lowercase{\endgroup\let~}\lycee@moins
\newcommand\lycee@moinsepais{\mathcode`\-="8000 }
\AtBeginDocument{\mathchardef\lycee@moinsorig=\mathcode`\-\relax}
\AtBeginDocument{\catcode`\_=12 \mathcode`\_="8000
  \catcode`\^=12 \mathcode`\^="8000 }
\newcommand\lycee@exposants{%
  \fontdimen13\textfont2=.48\fontdimen6\textfont2
  \fontdimen14\textfont2=.48\fontdimen6\textfont2
  \fontdimen15\textfont2=.48\fontdimen6\textfont2
  \fontdimen13\scriptfont2=.48\fontdimen6\scriptfont2
  \fontdimen14\scriptfont2=.48\fontdimen6\scriptfont2
  \fontdimen15\scriptfont2=.48\fontdimen6\scriptfont2}
\every@math@size\expandafter{\the\every@math@size\lycee@exposants}
\newlength{\retraitcalculs}
\setlength{\retraitcalculs}{2em}

\RequirePackage{tikz}
\usetikzlibrary{arrows.meta,calc,babel,shapes.misc}
\RequirePackage[most]{tcolorbox}
\tcbuselibrary{skins,breakable}

\definecolor{coulDef}{HTML}{1F4E79}
\definecolor{coulProp}{HTML}{7030A0}
\definecolor{coulRem}{HTML}{C55A11}
\definecolor{coulEx}{HTML}{2E7D32}
\definecolor{coulExo}{HTML}{B03A2E}
\definecolor{coulPart}{HTML}{1F4E79}
\definecolor{coulMeth}{HTML}{1B6E4F}
\definecolor{coulCourbe}{HTML}{0B5ED7}
\definecolor{coulCourbeB}{HTML}{C00000}
\definecolor{coulCourbeC}{HTML}{2E7D32}
\definecolor{coulGrille}{HTML}{8C8C8C}
\definecolor{coulRep}{HTML}{1B6E4F}
\definecolor{coulBar}{HTML}{2E7D32}

\newcommand{\R}{\mathbb{R}}
\newcommand{\Rs}{\mathbb{R}^{*}}
\renewcommand{\le}{\leqslant}
\renewcommand{\ge}{\geqslant}
\renewcommand{\approx}{\simeq}
\newcommand{\vect}[1]{\overrightarrow{#1}}
\newcommand{\norme}[1]{\left\lVert #1 \right\rVert}
\newcommand{\intervalle}[2]{\left[#1\,;#2\right]}
\RequirePackage{cancel}
\renewcommand{\CancelColor}{\color{coulExo}}
\newcommand{\fois}[1]{\times{\color{coulExo}#1}}
\newcommand{\barre}[1]{\cancel{#1}}

\tikzset{grille/.style={coulGrille,line width=0.4pt}}
\newcommand{\quadrillage}[4]{\draw[grille,step=1] (#1,#2) grid (#3,#4);}
\tikzset{
  croix/.style={cross out,draw,line width=0.6pt,inner sep=0pt,minimum size=4pt},
  croixdroite/.style={croix,rotate=45,line width=0.8pt},
}
\newcommand{\pt}[3]{\path (#1) node[croix]{} node[#2,font=\small,inner sep=2.5pt]{$#3$};}

\newcommand{\lycee@repere}[4]{%
  \draw[grille,step=1] (#1,#2) grid (#3,#4);
  \draw[-{Stealth[length=2mm]},thick] (#1,0)--({#3+0.4},0) node[below left=-1pt]{$x$};
  \draw[-{Stealth[length=2mm]},thick] (0,#2)--(0,{#4+0.4}) node[below left=-1pt]{$y$};
  \node[below left,font=\scriptsize,inner sep=1.5pt] at (0,0) {$0$};}
\newcommand{\lycee@gradx}[1]{\foreach \k/\l in {#1}{%
  \draw (\k,0.07)--(\k,-0.07) node[below,font=\scriptsize,inner sep=1.5pt]{$\l$};}}
\newcommand{\lycee@grady}[1]{\foreach \k/\l in {#1}{%
  \draw (0.07,\k)--(-0.07,\k) node[left,font=\scriptsize,inner sep=1.5pt]{$\l$};}}
\newcommand{\lycee@intff}[2]{\left[#1\,;#2\right]}
\newcommand{\lycee@intfo}[2]{\left[#1\,;#2\right[}
\newcommand{\lycee@intof}[2]{\left]#1\,;#2\right]}
\newcommand{\lycee@intoo}[2]{\left]#1\,;#2\right[}
\newcommand{\lycee@ens}[1]{\left\{#1\right\}}
\AtBeginDocument{%
  \providecommand{\repere}{\lycee@repere}%
  \providecommand{\gradx}{\lycee@gradx}%
  \providecommand{\grady}{\lycee@grady}%
  \providecommand{\Cf}{\mathcal{C}\sb{\scriptscriptstyle f}}%
  \providecommand{\Cg}{\mathcal{C}\sb{\scriptscriptstyle g}}%
  \providecommand{\intff}{\lycee@intff}%
  \providecommand{\intfo}{\lycee@intfo}%
  \providecommand{\intof}{\lycee@intof}%
  \providecommand{\intoo}{\lycee@intoo}%
  \providecommand{\ens}{\lycee@ens}}

\newlist{questions}{enumerate}{3}
\setlist[questions,1]{label=\arabic*\textdegree),leftmargin=*,itemsep=2pt,topsep=3pt}
\setlist[questions,2]{label=\alph*),leftmargin=*,itemsep=1pt,topsep=2pt}
\setlist[questions,3]{label=\roman*),leftmargin=*,itemsep=1pt,topsep=1pt}
\setlist[itemize]{label=\textbullet,leftmargin=*,itemsep=2pt,topsep=3pt}

\newcommand{\besoinplace}[1]{\par\penalty\@M\begingroup
  \dimen@=#1\relax\dimen@ii=\pagegoal\advance\dimen@ii-\pagetotal
  \ifdim\dimen@>\dimen@ii\ifdim\dimen@ii>\z@\vfil\fi\break\fi\endgroup}

\newcommand{\lycee@pied}{\thepage/\pageref{LastPage}}
\newcommand{\entete}[2]{%
  \pagestyle{fancy}\fancyhf{}%
  \fancyhead[L]{\footnotesize\color{coulPart}\textbf{#1}}%
  \fancyhead[R]{\footnotesize\color{coulPart}\textbf{#2}}%
  \fancyfoot[C]{\footnotesize\lycee@pied}%
  \renewcommand{\headrulewidth}{0.4pt}%
  \renewcommand{\headrule}{\hbox to\headwidth{%
    \color{coulPart!40}\leaders\hrule height \headrulewidth\hfill}}}

\RequirePackage{ccicons}
\newcommand{\lycee@licence}{{\scriptsize\color{gray}%
  \copyright\ Luc Destombe --- \ccbyncsa\ CC BY-NC-SA 4.0}}
\newif\iflycee@licence \lycee@licencetrue
\newcommand{\sanslicence}{\lycee@licencefalse}
\AtBeginDocument{\iflycee@licence\AddToHookNext{shipout/foreground}{%
  \put(0,0){\rlap{\kern\dimexpr1in+\hoffset+\oddsidemargin+\textwidth\relax
    \llap{\raisebox{-\dimexpr1in+\voffset+\topmargin+\headheight+\headsep
      +\textheight+\footskip\relax}{\lycee@licence}}}}}\fi}

\newcommand{\formulesserrees}{%
  \setlength{\abovedisplayskip}{3pt}\setlength{\belowdisplayskip}{3pt}%
  \setlength{\abovedisplayshortskip}{2pt}\setlength{\belowdisplayshortskip}{3pt}}

\setlength{\parindent}{0pt}

\RequirePackage[implicit=false,hidelinks,pdfencoding=auto,psdextra,bookmarksopen,
  bookmarksopenlevel=1,pdfcreator={},pdfproducer={}]{hyperref}
\RequirePackage{bookmark}
\hypersetup{pdfauthor={Luc Destombe},pdfkeywords={CC BY-NC-SA 4.0}}
\newcounter{lycee@signet}
\newcommand{\signet}[3][]{\stepcounter{lycee@signet}%
  \edef\lycee@dest{\if\relax\detokenize{#1}\relax
    signet.\arabic{lycee@signet}\else#1\fi}%
  \hypertarget{\lycee@dest}{}\bookmark[level=#2,dest=\lycee@dest]{#3}}
\pdfstringdefDisableCommands{%
  \def\R{R}\def\vect#1{#1}\def\Cf{Cf}\def\Cg{Cg}\def\fois#1{×#1}%
  \def\barre#1{#1}\def\intervalle#1#2{[#1;#2]}\def\intff#1#2{[#1;#2]}%
  \def\textsuperscript#1{#1}\def\dfrac#1#2{#1/#2}\def\frac#1#2{#1/#2}%
  \def\vec#1{#1⃗}}%
\RequirePackage[official]{eurosym}

\tcbset{exercice corrige/.style={
  enhanced,breakable,before upper=\raggedright,
  colback=white,colframe=coulExo,
  boxrule=0.8pt,arc=2pt,
  left=6pt,right=6pt,top=8pt,bottom=5pt,
  before skip=9pt,after skip=4pt,
  coltitle=white,fonttitle=\bfseries\small,
  attach boxed title to top left={xshift=8pt,yshift=-\tcboxedtitleheight/2},
  boxed title style={colback=coulExo,arc=2pt,boxrule=0pt}}}
\newcounter{exo}
\newtcolorbox{exercice}[1][]{exercice corrige,
  phantom={\signet[exercice-\arabic{exo}]{2}{Exercice \arabic{exo}}},
  title={Exercice \theexo\ifx\relax#1\relax\else{} --- #1\fi}}
\newenvironment{exo}[1][\relax]{\refstepcounter{exo}\begin{exercice}[#1]}{\end{exercice}}

\geometry{top=1.8cm,bottom=1cm,left=1cm,right=1cm,headheight=14pt,footskip=0.6cm}
\renewcommand{\lycee@pied}{\thepage}

\renewcommand{\theexo}{\ifnum\value{exo}<10 0\fi\arabic{exo}}
\tcbset{exercice corrige/.style={
  enhanced,breakable,before upper=\raggedright,
  colback=white,colframe=coulExo,
  boxrule=0.9pt,arc=2pt,
  left=8pt,right=8pt,top=8pt,bottom=6pt,
  before skip=8pt,after skip=6pt,
  coltitle=white,fonttitle=\bfseries,
  attach boxed title to top left={xshift=8pt,yshift=-\tcboxedtitleheight/2},
  boxed title style={colback=coulExo,arc=2pt,boxrule=0pt}}}

\newcommand{\entetefiche}[2]{%
  \begin{center}
    \begin{tcolorbox}[enhanced,width=0.92\linewidth,colback=white,colframe=coulPart,
        boxrule=1.4pt,arc=3pt,halign=center,top=6pt,bottom=6pt]
      {\large\bfseries\color{coulPart} CORRIGÉ --- FICHE D'EXERCICES}\\[3pt]
      {\normalsize\bfseries #1}\\[2pt]
      {\normalsize #2}
    \end{tcolorbox}
  \end{center}}

\newcounter{partie}
\newcommand{\partie}[1]{%
  \stepcounter{partie}%
  \besoinplace{8\baselineskip}\vspace{4pt}%
  \begin{tcolorbox}[enhanced,colback=coulPart!8,colframe=coulPart,boxrule=0pt,
      leftrule=3pt,arc=0pt,left=6pt,right=4pt,top=3pt,bottom=3pt,
      before skip=6pt,after skip=2pt]
    \bfseries\color{coulPart}\Roman{partie} --- #1\signet{1}{\Roman{partie} --- #1}
  \end{tcolorbox}}
\newcommand{\partiesansnum}[1]{%
  \besoinplace{8\baselineskip}\vspace{4pt}%
  \begin{tcolorbox}[enhanced,colback=coulPart!8,colframe=coulPart,boxrule=0pt,
      leftrule=3pt,arc=0pt,left=6pt,right=4pt,top=3pt,bottom=3pt,
      before skip=6pt,after skip=2pt]
    \bfseries\color{coulPart}#1\signet{1}{#1}
  \end{tcolorbox}}

\setlist[questions,1]{label=\arabic*\textdegree),leftmargin=*,itemsep=3pt,topsep=3pt}
\setlist[questions,2]{label=\alph*),leftmargin=*,itemsep=2pt,topsep=2pt}
\setlist[questions,3]{label=\roman*),leftmargin=*,itemsep=1pt,topsep=1pt}
\newlist{expr}{enumerate}{1}
\setlist[expr]{label={\Alph*\ $=$},leftmargin=*,itemsep=2pt,topsep=3pt}

\newcommand{\res}[1]{{\color{coulRep}#1}}
\newcommand{\rem}[1]{\par\smallskip{\small\itshape\color{black!60}$\rightarrow$ #1}}
\newenvironment{calculs}{\par\smallskip\noindent$\displaystyle\begin{aligned}[t]}%
  {\end{aligned}$\par\smallskip}
\newcommand{\justif}[1]{\quad\text{\small\itshape\color{black!60}#1}}
\newcommand{\cas}[1]{\par\smallskip\noindent\textbf{\color{coulExo}#1}\nobreak}
\newcommand{\explique}[2]{\par\smallskip\noindent
  \begin{minipage}[t]{0.6\linewidth}\raggedright #1\end{minipage}\hfill
  \begin{minipage}[t]{0.37\linewidth}\raggedright\leavevmode
    {\small\itshape\color{black!60}#2}\end{minipage}\par}

\setlength{\parskip}{2pt}

\makeatother
\geometry{top=1.8cm,bottom=1.6cm,left=1.8cm,right=1.8cm,headheight=14pt,footskip=30pt}
\entete{Chapitre 3 --- Le second degré --- corrigé}{Première spécialité}
\setlength{\parindent}{0pt}
\setlength{\jot}{4pt}

\let\casligne\cas
\renewcommand{\cas}[1]{\casligne{#1}\par\nopagebreak}

\usepackage{tkz-tab}
\usepackage{pgfplots}
\pgfplotsset{compat=1.18}
\pgfplotsset{
  repere/.style={
    axis lines=middle,
    axis line style={-{Stealth[length=2mm]},thick},
    label style={font=\scriptsize},
    tick label style={font=\scriptsize},
    grid=both,
    major grid style={grille},
    minor tick num=0,
    tick align=inside,
    clip mode=individual,
    xlabel={$x$},ylabel={$y$},
  },
  courbe/.style={coulCourbe,line width=1pt,samples=100},
  croix/.style={only marks,mark=x,mark size=2pt,mark options={line width=0.6pt},black},
}

\begin{document}

\entetefiche{Chapitre 3 --- Le second degré}{Première spécialité mathématiques}

\partie{Forme canonique}

\begin{exo}[Lire, écrire et développer une forme canonique]
\explique{On compare chaque expression à la forme du cours $a(x-\alpha)^{2}+\beta$ :
$a$ est le nombre devant la parenthèse, $\alpha$ le nombre \emph{retranché} à $x$,
$\beta$ le nombre ajouté à la fin.}{avant de lire, on réécrit l'expression exactement
sous cette forme : « $x-$ » dans la parenthèse, « $+$ » devant $\beta$}

\cas{1) a)}
\begin{calculs}
f(x) &= (x-3)^{2}+5 &&\justif{déjà sous la forme $a(x-\alpha)^{2}+\beta$}
\end{calculs}
\explique{$a=\res{1}$, $\alpha=\res{3}$ et $\beta=\res{5}$.}{rien devant la
parenthèse : $a=1$}

\cas{1) b)}
\begin{calculs}
g(x) &= (x+2)^{2}-7 &&\justif{on recopie}\\
g(x) &= \bigl(x-(-2)\bigr)^{2}+(-7) &&\justif{$x+2=x-(-2)$ ; $-7=+(-7)$}
\end{calculs}
\explique{$a=\res{1}$, $\alpha=\res{-2}$ et $\beta=\res{-7}$.}{le signe de $\alpha$ est
le contraire de celui qu'on lit dans la parenthèse}

\cas{1) c)}
\begin{calculs}
h(x) &= 2(x+1)^{2}+4 &&\justif{on recopie}\\
h(x) &= 2\bigl(x-(-1)\bigr)^{2}+4 &&\justif{$x+1=x-(-1)$}
\end{calculs}
\explique{$a=\res{2}$, $\alpha=\res{-1}$ et $\beta=\res{4}$.}{}

\cas{1) d)}
\begin{calculs}
k(x) &= -3(x-4)^{2}-1 &&\justif{on recopie}\\
k(x) &= -3(x-4)^{2}+(-1) &&\justif{$\beta$ garde son signe : $-1$}
\end{calculs}
\explique{$a=\res{-3}$, $\alpha=\res{4}$ et $\beta=\res{-1}$.}{le signe « $-$ » devant
le $3$ appartient à $a$}

\cas{1) e)}
\begin{calculs}
m(x) &= (x+5)^{2} &&\justif{on recopie}\\
m(x) &= \bigl(x-(-5)\bigr)^{2}+0 &&\justif{rien n'est ajouté : $\beta=0$}
\end{calculs}
\explique{$a=\res{1}$, $\alpha=\res{-5}$ et $\beta=\res{0}$.}{}

\cas{1) f)}
\begin{calculs}
n(x) &= -x^{2}+6 &&\justif{on recopie}\\
n(x) &= -(x-0)^{2}+6 &&\justif{$x^{2}=(x-0)^{2}$ ; le « $-$ » devant donne $a=-1$}
\end{calculs}
\explique{$a=\res{-1}$, $\alpha=\res{0}$ et $\beta=\res{6}$.}{}
\rem{Erreur fréquente : lire $\alpha=2$ dans $(x+2)^{2}$. $\alpha$ est le nombre que
l'on \emph{retranche} à $x$ : dans $(x+2)$, $\alpha=-2$.}

\cas{2) a)}
\explique{On remplace $a$, $\alpha$ et $\beta$ dans $P(x)=a(x-\alpha)^{2}+\beta$, puis on
simplifie les signes.}{même travail que la question 1, dans l'autre sens}
\begin{calculs}
P(x) &= 1\times(x-2)^{2}+(-3) &&\justif{$a=1$, $\alpha=2$ et $\beta=-3$}\\
P(x) &= \res{(x-2)^{2}-3} &&\justif{$1\times$ ne s'écrit pas ; $+(-3)=-3$}
\end{calculs}
\cas{2) b)}
\begin{calculs}
P(x) &= -2\bigl(x-(-1)\bigr)^{2}+5 &&\justif{$a=-2$, $\alpha=-1$ et $\beta=5$}\\
P(x) &= \res{-2(x+1)^{2}+5} &&\justif{$x-(-1)=x+1$}
\end{calculs}
\cas{2) c)}
\begin{calculs}
P(x) &= 3(x-4)^{2}+0 &&\justif{$a=3$, $\alpha=4$ et $\beta=0$}\\
P(x) &= \res{3(x-4)^{2}} &&\justif{on n'écrit pas $+0$}
\end{calculs}
\cas{2) d)}
\begin{calculs}
P(x) &= 5(x-0)^{2}+(-1) &&\justif{$a=5$, $\alpha=0$ et $\beta=-1$}\\
P(x) &= \res{5x^{2}-1} &&\justif{$(x-0)^{2}=x^{2}$}
\end{calculs}
\rem{Contrôle : relire la forme obtenue comme à la question 1 ; on doit retrouver les
$a$, $\alpha$ et $\beta$ de l'énoncé.}

\cas{3)}
\explique{On développe d'abord le carré (identité remarquable, gardée entre
parenthèses), puis on multiplie par $a$, et enfin on réduit.}{l'ordre compte : le carré
avant la multiplication par $a$}
\begin{calculs}
f(x) &= (x-3)^{2}+5 &&\justif{on recopie}\\
f(x) &= x^{2}-6x+9+5 &&\justif{$(u-v)^{2}=u^{2}-2uv+v^{2}$, avec $u=x$ et $v=3$}\\
f(x) &= \res{x^{2}-6x+14} &&\justif{on réduit}
\end{calculs}
\begin{calculs}
g(x) &= (x+2)^{2}-7 &&\justif{on recopie}\\
g(x) &= x^{2}+4x+4-7 &&\justif{$(u+v)^{2}=u^{2}+2uv+v^{2}$}\\
g(x) &= \res{x^{2}+4x-3} &&\justif{on réduit}
\end{calculs}
\begin{calculs}
h(x) &= 2(x+1)^{2}+4 &&\justif{on recopie}\\
h(x) &= 2\left(x^{2}+2x+1\right)+4 &&\justif{le carré développé, entre parenthèses}\\
h(x) &= 2x^{2}+4x+2+4 &&\justif{on distribue le $2$}\\
h(x) &= \res{2x^{2}+4x+6} &&\justif{on réduit}
\end{calculs}
\begin{calculs}
k(x) &= -3(x-4)^{2}-1 &&\justif{on recopie}\\
k(x) &= -3\left(x^{2}-8x+16\right)-1 &&\justif{le carré développé, entre parenthèses}\\
k(x) &= -3x^{2}+24x-48-1 &&\justif{on distribue le $-3$ : tous les signes changent}\\
k(x) &= \res{-3x^{2}+24x-49} &&\justif{on réduit}
\end{calculs}
\rem{Erreur fréquente : faire entrer $a$ dans le carré, $2(x+1)^{2}=(2x+2)^{2}$. C'est
faux : $(2x+2)^{2}=4(x+1)^{2}$. Contrôle rapide : en $x=0$, la forme canonique de $k$
donne $-3\times 16-1=-49$, qui est bien le terme constant trouvé.}
\end{exo}

\begin{exo}[Compléter un carré ($a=1$)]
\explique{$x^{2}+kx$ est le début du carré $\left(x+\dfrac{k}{2}\right)^{2}$ : on ajoute
le terme qui manque, $\left(\dfrac{k}{2}\right)^{2}$, et on le retire aussitôt.}{méthode
1 : la valeur de l'expression ne change pas}

\cas{A}
\begin{calculs}
A &= x^{2}+8x-3 &&\justif{$x^{2}+8x$ est le début de $(x+4)^{2}=x^{2}+8x+16$}\\
A &= x^{2}+8x+16-16-3 &&\justif{on ajoute $16$ et on le retire}\\
A &= \res{(x+4)^{2}-19} &&\justif{on regroupe le carré ; $-16-3=-19$}
\end{calculs}
\cas{B}
\begin{calculs}
B &= x^{2}-6x+2 &&\justif{$x^{2}-6x$ est le début de $(x-3)^{2}=x^{2}-6x+9$}\\
B &= x^{2}-6x+9-9+2 &&\justif{on ajoute $9$ et on le retire}\\
B &= \res{(x-3)^{2}-7} &&\justif{$-9+2=-7$}
\end{calculs}
\cas{C}
\begin{calculs}
C &= x^{2}-3x+1 &&\justif{moitié de $-3$ : $-\frac{3}{2}$ ; $\left(x-\frac{3}{2}\right)^{2}=x^{2}-3x+\frac{9}{4}$}\\
C &= x^{2}-3x+\dfrac{9}{4}-\dfrac{9}{4}+1 &&\justif{on ajoute $\frac{9}{4}$ et on le retire}\\
C &= \left(x-\dfrac{3}{2}\right)^{2}-\dfrac{9}{4}+\dfrac{1\fois{4}}{1\fois{4}} &&\justif{même dénominateur $4$}\\
C &= \res{\left(x-\dfrac{3}{2}\right)^{2}-\dfrac{5}{4}} &&\justif{$-9+4=-5$}
\end{calculs}
\cas{D}
\begin{calculs}
D &= x^{2}+5x+2 &&\justif{moitié de $5$ : $\frac{5}{2}$ ; $\left(x+\frac{5}{2}\right)^{2}=x^{2}+5x+\frac{25}{4}$}\\
D &= x^{2}+5x+\dfrac{25}{4}-\dfrac{25}{4}+2 &&\justif{on ajoute $\frac{25}{4}$ et on le retire}\\
D &= \left(x+\dfrac{5}{2}\right)^{2}-\dfrac{25}{4}+\dfrac{2\fois{4}}{1\fois{4}} &&\justif{même dénominateur $4$}\\
D &= \res{\left(x+\dfrac{5}{2}\right)^{2}-\dfrac{17}{4}} &&\justif{$-25+8=-17$}
\end{calculs}
\cas{E}
\begin{calculs}
E &= x^{2}-12x+35 &&\justif{$x^{2}-12x$ est le début de $(x-6)^{2}=x^{2}-12x+36$}\\
E &= x^{2}-12x+36-36+35 &&\justif{on ajoute $36$ et on le retire}\\
E &= \res{(x-6)^{2}-1} &&\justif{$-36+35=-1$}
\end{calculs}
\cas{F}
\begin{calculs}
F &= x^{2}-9x &&\justif{moitié de $-9$ : $-\frac{9}{2}$ ; $\left(x-\frac{9}{2}\right)^{2}=x^{2}-9x+\frac{81}{4}$}\\
F &= x^{2}-9x+\dfrac{81}{4}-\dfrac{81}{4} &&\justif{on ajoute $\frac{81}{4}$ et on le retire}\\
F &= \res{\left(x-\dfrac{9}{2}\right)^{2}-\dfrac{81}{4}} &&\justif{pas de terme constant : rien d'autre à réduire}
\end{calculs}
\rem{Dans la parenthèse : la \emph{moitié} du coefficient de $x$, avec son signe. On
retire le \emph{carré} de cette moitié ($16$ pour $8x$), pas le coefficient lui-même.
Contrôle : on développe la forme trouvée, par exemple $(x+4)^{2}-19=x^{2}+8x+16-19$,
soit $x^{2}+8x-3$.}
\end{exo}

\begin{exo}[Compléter un carré ($a\ne 1$)]
\explique{Quand $a\ne 1$, on met d'abord $a$ en facteur dans les deux premiers termes :
la parenthèse commence alors par $x^{2}$, et on complète le carré comme à l'exercice
2.}{méthode 1, exemples b) et c)}

\cas{A}
\begin{calculs}
A &= 2x^{2}+8x+3 &&\justif{on recopie}\\
A &= 2\left(x^{2}+4x\right)+3 &&\justif{$2$ en facteur dans les deux premiers termes : $8x=2\times 4x$}\\
A &= 2\left[(x+2)^{2}-4\right]+3 &&\justif{$x^{2}+4x$ est le début de $(x+2)^{2}=x^{2}+4x+4$}\\
A &= 2(x+2)^{2}-8+3 &&\justif{on distribue le $2$ : $2\times(-4)=-8$}\\
A &= \res{2(x+2)^{2}-5} &&\justif{on réduit}
\end{calculs}
\cas{B}
\begin{calculs}
B &= -x^{2}+3x+2 &&\justif{on recopie}\\
B &= -\left(x^{2}-3x\right)+2 &&\justif{$-1$ en facteur : les signes changent dans la parenthèse}\\
B &= -\left[\left(x-\dfrac{3}{2}\right)^{2}-\dfrac{9}{4}\right]+2 &&\justif{$x^{2}-3x$ est le début de $\left(x-\frac{3}{2}\right)^{2}$}\\
B &= -\left(x-\dfrac{3}{2}\right)^{2}+\dfrac{9}{4}+\dfrac{2\fois{4}}{1\fois{4}} &&\justif{on distribue le « $-$ » : $-\left(-\frac{9}{4}\right)=+\frac{9}{4}$}\\
B &= \res{-\left(x-\dfrac{3}{2}\right)^{2}+\dfrac{17}{4}} &&\justif{$9+8=17$}
\end{calculs}
\cas{C}
\begin{calculs}
C &= 4x^{2}-8x+4 &&\justif{on recopie}\\
C &= 4\left(x^{2}-2x\right)+4 &&\justif{$4$ en facteur dans les deux premiers termes}\\
C &= 4\left[(x-1)^{2}-1\right]+4 &&\justif{$x^{2}-2x$ est le début de $(x-1)^{2}=x^{2}-2x+1$}\\
C &= 4(x-1)^{2}-4+4 &&\justif{on distribue le $4$}\\
C &= \res{4(x-1)^{2}} &&\justif{ici $\beta=0$}
\end{calculs}
\rem{On pouvait voir directement $4x^{2}-8x+4=4\left(x^{2}-2x+1\right)$ : la parenthèse
est déjà l'identité remarquable $(x-1)^{2}$.}
\cas{D}
\begin{calculs}
D &= -x^{2}+5x-4 &&\justif{on recopie}\\
D &= -\left(x^{2}-5x\right)-4 &&\justif{$-1$ en facteur : les signes changent}\\
D &= -\left[\left(x-\dfrac{5}{2}\right)^{2}-\dfrac{25}{4}\right]-4 &&\justif{$x^{2}-5x$ est le début de $\left(x-\frac{5}{2}\right)^{2}$}\\
D &= -\left(x-\dfrac{5}{2}\right)^{2}+\dfrac{25}{4}-\dfrac{4\fois{4}}{1\fois{4}} &&\justif{on distribue le « $-$ » ; même dénominateur}\\
D &= \res{-\left(x-\dfrac{5}{2}\right)^{2}+\dfrac{9}{4}} &&\justif{$25-16=9$}
\end{calculs}
\cas{E}
\begin{calculs}
E &= 2x^{2}+2x-1 &&\justif{on recopie}\\
E &= 2\left(x^{2}+x\right)-1 &&\justif{$2$ en facteur : $2x=2\times x$}\\
E &= 2\left[\left(x+\dfrac{1}{2}\right)^{2}-\dfrac{1}{4}\right]-1 &&\justif{$x^{2}+x$ est le début de $\left(x+\frac{1}{2}\right)^{2}=x^{2}+x+\frac{1}{4}$}\\
E &= 2\left(x+\dfrac{1}{2}\right)^{2}-\dfrac{1}{2}-\dfrac{1\fois{2}}{1\fois{2}} &&\justif{$2\times\left(-\frac{1}{4}\right)=-\frac{1}{2}$ ; même dénominateur}\\
E &= \res{2\left(x+\dfrac{1}{2}\right)^{2}-\dfrac{3}{2}} &&\justif{$-1-2=-3$}
\end{calculs}
\cas{F}
\begin{calculs}
F &= -2x^{2}+8x-5 &&\justif{on recopie}\\
F &= -2\left(x^{2}-4x\right)-5 &&\justif{$-2$ en facteur : $8x=-2\times(-4x)$}\\
F &= -2\left[(x-2)^{2}-4\right]-5 &&\justif{$x^{2}-4x$ est le début de $(x-2)^{2}$}\\
F &= -2(x-2)^{2}+8-5 &&\justif{on distribue le $-2$ : $-2\times(-4)=8$}\\
F &= \res{-2(x-2)^{2}+3} &&\justif{on réduit}
\end{calculs}
\rem{Erreurs fréquentes : oublier de multiplier par $a$ le terme retiré (en A,
$2\times(-4)=-8$, et non $-4$) ; oublier de changer les signes quand on met un nombre
négatif en facteur. Contrôle : on développe le résultat ; en F,
$-2\left(x^{2}-4x+4\right)+3=-2x^{2}+8x-5$.}
\end{exo}

\partie{Calcul de $\alpha$ et $\beta$, factorisation}

\begin{exo}[Forme canonique par $\alpha$ et $\beta$]
\explique{On lit $a$, $b$ et $c$, puis $\alpha=-\dfrac{b}{2a}$ et $\beta=P(\alpha)$ :
en $x=\alpha$ le carré est nul, il ne reste que $\beta$.}{théorème 1 (complet) : plus
rapide que de compléter le carré}

\cas{1)}
\explique{$P(x)=x^{2}-6x+1$ : $a=1$, $b=-6$ et $c=1$.}{on lit les coefficients avec
leur signe}
\begin{calculs}
\alpha &= -\dfrac{-6}{2\times 1}=3 &&\justif{$\alpha=-\frac{b}{2a}$ : deux signes « $-$ » donnent $+$}\\
\beta &= P(3)=9-18+1 &&\justif{$\beta=P(\alpha)$ : $3^{2}-6\times 3+1$}\\
\beta &= -8 &&
\end{calculs}
\explique{Forme canonique : $P(x)=\res{(x-3)^{2}-8}$.}{$a(x-\alpha)^{2}+\beta$ avec
$a=1$}
\begin{calculs}
(x-3)^{2}-8 &= x^{2}-6x+9-8 &&\justif{vérification : on développe}\\
(x-3)^{2}-8 &= x^{2}-6x+1 &&\justif{on retrouve $P(x)$}
\end{calculs}

\cas{2)}
\explique{$Q(x)=3x^{2}+6x-2$ : $a=3$, $b=6$ et $c=-2$.}{}
\begin{calculs}
\alpha &= -\dfrac{6}{2\times 3}=-1 &&\justif{$\alpha=-\frac{b}{2a}$}\\
\beta &= Q(-1)=3-6-2 &&\justif{$3\times(-1)^{2}+6\times(-1)-2$}\\
\beta &= -5 &&
\end{calculs}
\explique{Forme canonique : $Q(x)=3\bigl(x-(-1)\bigr)^{2}-5$, soit
$Q(x)=\res{3(x+1)^{2}-5}$.}{devant le carré, on écrit $a=3$}
\begin{calculs}
3(x+1)^{2}-5 &= 3\left(x^{2}+2x+1\right)-5 &&\justif{vérification : on développe le carré}\\
3(x+1)^{2}-5 &= 3x^{2}+6x+3-5 &&\justif{on distribue le $3$}\\
3(x+1)^{2}-5 &= 3x^{2}+6x-2 &&\justif{on retrouve $Q(x)$}
\end{calculs}

\cas{3)}
\explique{$R(x)=-x^{2}-4x-9$ : $a=-1$, $b=-4$ et $c=-9$.}{trois coefficients négatifs :
attention aux signes}
\begin{calculs}
\alpha &= -\dfrac{-4}{2\times(-1)}=-2 &&\justif{$-\frac{-4}{-2}$ : trois signes « $-$ », le résultat est négatif}\\
\beta &= R(-2)=-4+8-9 &&\justif{$-(-2)^{2}=-4$ : le carré d'abord, puis le « $-$ »}\\
\beta &= -5 &&
\end{calculs}
\explique{Forme canonique : $R(x)=\res{-(x+2)^{2}-5}$.}{$a=-1$ s'écrit « $-$ » devant le
carré}
\begin{calculs}
-(x+2)^{2}-5 &= -\left(x^{2}+4x+4\right)-5 &&\justif{vérification}\\
-(x+2)^{2}-5 &= -x^{2}-4x-9 &&\justif{on retrouve $R(x)$}
\end{calculs}
\rem{Devant le carré, on écrit toujours $a$ : $3$ en 2), $-1$ en 3). L'oublier est
l'erreur la plus fréquente ; la vérification par le développement la détecte aussitôt.}
\end{exo}

\begin{exo}[De la forme canonique à la factorisation]
\explique{Dans $a(x-\alpha)^{2}+\beta$, on regarde les signes de $a$ et de $\beta$.
Signes contraires : on fait apparaître $u^{2}-v^{2}=(u-v)(u+v)$. Même signe : le
trinôme ne s'annule jamais, il ne se factorise pas.}{méthode 2}

\cas{A}
\explique{$A=x^{2}-4x-5$ : $a=1$, $b=-4$ et $c=-5$.}{}
\begin{calculs}
\alpha &= -\dfrac{-4}{2}=2 &&\justif{$\alpha=-\frac{b}{2a}$}\\
\beta &= 4-8-5 &&\justif{$\beta=A(2)=2^{2}-4\times 2-5$}\\
\beta &= -9 &&
\end{calculs}
\begin{calculs}
A &= (x-2)^{2}-9 &&\justif{forme canonique : $a=1>0$ et $\beta<0$, signes contraires}\\
A &= (x-2)^{2}-3^{2} &&\justif{forme $u^{2}-v^{2}$ avec $u=x-2$ et $v=3$}\\
A &= (x-2-3)(x-2+3) &&\justif{$(u-v)(u+v)$}\\
A &= \res{(x-5)(x+1)} &&\justif{on réduit chaque facteur}
\end{calculs}
\explique{Contrôle : $(x-5)(x+1)=x^{2}+x-5x-5$, soit $x^{2}-4x-5$.}{on développe pour
retrouver l'énoncé}

\cas{B}
\explique{$B=x^{2}+4x+7$ : $a=1$, $b=4$ et $c=7$.}{}
\begin{calculs}
\alpha &= -\dfrac{4}{2}=-2 &&\\
\beta &= 4-8+7 &&\justif{$\beta=B(-2)=(-2)^{2}+4\times(-2)+7$}\\
\beta &= 3 &&\\
B &= \res{(x+2)^{2}+3} &&\justif{forme canonique : $a>0$ et $\beta>0$, même signe}
\end{calculs}
\explique{$(x+2)^{2}+3$ est une \emph{somme}, pas une différence de deux carrés. Un
carré est positif, donc $(x+2)^{2}+3\ge 3$ : $B$ ne vaut jamais $0$.}{un produit de
facteurs du premier degré, lui, s'annulerait}
\explique{Le trinôme $B$ \res{ne se factorise pas}.}{on répond à la question par une
phrase}

\cas{C}
\explique{$C=2x^{2}+x-3$ : $a=2$, $b=1$ et $c=-3$.}{}
\begin{calculs}
\alpha &= -\dfrac{1}{2\times 2}=-\dfrac{1}{4} &&\justif{$\alpha=-\frac{b}{2a}$}\\
\beta &= 2\times\dfrac{1}{16}-\dfrac{1}{4}-3 &&\justif{$\beta=C\left(-\frac{1}{4}\right)$ ; $\left(-\frac{1}{4}\right)^{2}=\frac{1}{16}$}\\
\beta &= \dfrac{1}{8}-\dfrac{1\fois{2}}{4\fois{2}}-\dfrac{3\fois{8}}{1\fois{8}} &&\justif{$\frac{2}{16}=\frac{1}{8}$ ; même dénominateur $8$}\\
\beta &= -\dfrac{25}{8} &&\justif{$1-2-24=-25$}
\end{calculs}
\begin{calculs}
C &= \res{2\left(x+\dfrac{1}{4}\right)^{2}-\dfrac{25}{8}} &&\justif{forme canonique : $a>0$ et $\beta<0$}\\
C &= 2\left[\left(x+\dfrac{1}{4}\right)^{2}-\dfrac{25}{16}\right] &&\justif{$2$ en facteur : $-\frac{25}{8}=2\times\left(-\frac{25}{16}\right)$}\\
C &= 2\left[\left(x+\dfrac{1}{4}\right)^{2}-\left(\dfrac{5}{4}\right)^{2}\right] &&\justif{$\frac{25}{16}=\left(\frac{5}{4}\right)^{2}$ : forme $u^{2}-v^{2}$}\\
C &= 2\left(x+\dfrac{1}{4}-\dfrac{5}{4}\right)\left(x+\dfrac{1}{4}+\dfrac{5}{4}\right) &&\justif{$(u-v)(u+v)$}\\
C &= 2(x-1)\left(x+\dfrac{3}{2}\right) &&\justif{$\frac{1-5}{4}=-1$ et $\frac{1+5}{4}=\frac{3}{2}$}\\
C &= \res{(x-1)(2x+3)} &&\justif{le $2$ entre dans la seconde parenthèse}
\end{calculs}
\rem{On met $a$ en facteur dans la forme canonique \emph{avant} d'utiliser $u^{2}-v^{2}$ :
le crochet commence alors par un carré seul. Contrôle :
$(x-1)(2x+3)=2x^{2}+3x-2x-3$, soit $2x^{2}+x-3$.}

\cas{D}
\explique{$D=-x^{2}-2x-3$ : $a=-1$, $b=-2$ et $c=-3$.}{}
\begin{calculs}
\alpha &= -\dfrac{-2}{2\times(-1)}=-1 &&\justif{$-\frac{-2}{-2}$ : résultat négatif}\\
\beta &= -1+2-3 &&\justif{$\beta=D(-1)=-(-1)^{2}-2\times(-1)-3$}\\
\beta &= -2 &&\\
D &= \res{-(x+1)^{2}-2} &&\justif{forme canonique : $a<0$ et $\beta<0$, même signe}\\
D &= -\left[(x+1)^{2}+2\right] &&\justif{$-1$ en facteur : les signes changent}
\end{calculs}
\explique{Le crochet est une somme, pas de la forme $u^{2}-v^{2}$ ; il vaut au moins
$2$, donc $D\le -2$ : $D$ ne s'annule jamais et \res{ne se factorise pas}.}{$a$ et
$\beta$ de même signe}
\end{exo}

\begin{exo}[Trois formes d'un même trinôme]
\cas{1)}
\explique{$f(x)=2x^{2}-12x+10$ : $a=2$, $b=-12$ et $c=10$.}{}
\begin{calculs}
\alpha &= -\dfrac{-12}{2\times 2}=3 &&\justif{$\alpha=-\frac{b}{2a}$}\\
\beta &= f(3)=18-36+10 &&\justif{$2\times 3^{2}-12\times 3+10$}\\
\beta &= -8 &&
\end{calculs}
\explique{Forme canonique : $f(x)=\res{2(x-3)^{2}-8}$.}{$a$ devant le carré}

\cas{2)}
\explique{$a=2>0$ et $\beta=-8<0$ : signes contraires, la forme canonique se
factorise.}{méthode 2}
\begin{calculs}
f(x) &= 2\left[(x-3)^{2}-4\right] &&\justif{$2$ en facteur : $-8=2\times(-4)$}\\
f(x) &= 2\left[(x-3)^{2}-2^{2}\right] &&\justif{forme $u^{2}-v^{2}$ dans le crochet}\\
f(x) &= 2(x-3-2)(x-3+2) &&\justif{$(u-v)(u+v)$}\\
f(x) &= \res{2(x-5)(x-1)} &&\justif{on réduit}
\end{calculs}

\cas{3)}
\explique{On développe le produit des deux parenthèses, puis on multiplie par
$2$.}{passage de la forme factorisée à la forme développée}
\begin{calculs}
f(x) &= 2(x-5)(x-1) &&\justif{on recopie}\\
f(x) &= 2\left(x^{2}-x-5x+5\right) &&\justif{double distributivité, entre parenthèses}\\
f(x) &= 2\left(x^{2}-6x+5\right) &&\justif{on réduit}\\
f(x) &= 2x^{2}-12x+10 &&\justif{on distribue le $2$}
\end{calculs}
\explique{On retrouve la forme développée de l'énoncé : les trois écritures sont bien
égales.}{}

\cas{4)}
\explique{Pour chaque nombre, on choisit la forme où le calcul est immédiat.}{}
\begin{calculs}
f(0) &= 2\times 0^{2}-12\times 0+10 &&\justif{forme développée : les termes en $x$ sont nuls}\\
f(0) &= \res{10} &&\justif{c'est le terme constant $c$}
\end{calculs}
\begin{calculs}
f(3) &= 2(3-3)^{2}-8 &&\justif{forme canonique : en $\alpha=3$, le carré est nul}\\
f(3) &= \res{-8} &&\justif{c'est $\beta$}
\end{calculs}
\begin{calculs}
f(1) &= 2(1-5)(1-1) &&\justif{forme factorisée : le facteur $(x-1)$ est nul}\\
f(1) &= \res{0} &&\justif{$1$ est une valeur qui annule $f$}
\end{calculs}
\rem{Chaque forme a son usage : la développée pour $f(0)$, la canonique pour
$f(\alpha)$, la factorisée pour les valeurs qui annulent $f$.}
\end{exo}

\partie{Équations du second degré}

\begin{exo}[Factoriser plutôt que calculer $\Delta$]
\explique{Quand on voit un facteur commun ou une identité remarquable, on factorise
directement, puis on utilise : un produit est nul si et seulement si l'un de ses
facteurs est nul.}{méthode 3, exemple d) : pas de discriminant}

\cas{1)}
\begin{calculs}
9x^{2}-4=0 &\iff (3x-2)(3x+2)=0 &&\justif{$u^{2}-v^{2}$ avec $u=3x$ et $v=2$}\\
&\iff 3x-2=0 \text{ ou } 3x+2=0 &&\justif{produit nul}\\
&\iff x=\dfrac{2}{3} \text{ ou } x=-\dfrac{2}{3} &&
\end{calculs}
\explique{$S=\res{\left\{-\dfrac{2}{3}\,;\dfrac{2}{3}\right\}}$.}{on range les
solutions dans l'ordre croissant}
\cas{2)}
\begin{calculs}
x^{2}-6x=0 &\iff x(x-6)=0 &&\justif{facteur commun $x$}\\
&\iff x=0 \text{ ou } x-6=0 &&\justif{produit nul}
\end{calculs}
\explique{$S=\res{\{0\,;6\}}$.}{}
\rem{Erreur fréquente : diviser par $x$ et écrire $x=6$ ; on perd la solution $0$.}
\cas{3)}
\begin{calculs}
2x^{2}-18=0 &\iff 2\left(x^{2}-9\right)=0 &&\justif{facteur commun $2$}\\
&\iff 2(x-3)(x+3)=0 &&\justif{$x^{2}-9=x^{2}-3^{2}$}\\
&\iff x=3 \text{ ou } x=-3 &&\justif{le facteur $2$ ne s'annule pas}
\end{calculs}
\explique{$S=\res{\{-3\,;3\}}$.}{}
\cas{4)}
\begin{calculs}
x^{2}-5=0 &\iff \left(x-\sqrt{5}\right)\left(x+\sqrt{5}\right)=0 &&\justif{$5=\left(\sqrt{5}\right)^{2}$ : forme $u^{2}-v^{2}$}\\
&\iff x=\sqrt{5} \text{ ou } x=-\sqrt{5} &&\justif{produit nul}
\end{calculs}
\explique{$S=\res{\left\{-\sqrt{5}\,;\sqrt{5}\right\}}$.}{}
\cas{5)}
\begin{calculs}
(3x-2)(x+5)=0 &\iff 3x-2=0 \text{ ou } x+5=0 &&\justif{déjà un produit : on ne développe surtout pas}\\
&\iff x=\dfrac{2}{3} \text{ ou } x=-5 &&
\end{calculs}
\explique{$S=\res{\left\{-5\,;\dfrac{2}{3}\right\}}$.}{}
\cas{6)}
\begin{calculs}
x^{2}-10x+25=0 &\iff (x-5)^{2}=0 &&\justif{$u^{2}-2uv+v^{2}$ avec $u=x$ et $v=5$}\\
&\iff x-5=0 &&\justif{un carré est nul si et seulement si le nombre est nul}
\end{calculs}
\explique{$S=\res{\{5\}}$.}{une seule solution}
\rem{Facteur commun, identité remarquable, produit déjà factorisé : calculer $\Delta$
donnerait le même résultat, mais ferait perdre du temps.}
\end{exo}

\begin{exo}[Avec le discriminant]
\explique{On écrit l'équation sous la forme $ax^{2}+bx+c=0$, on lit $a$, $b$, $c$, on
calcule $\Delta=b^{2}-4ac$, et son signe donne le nombre de solutions.}{méthode 3 et
théorème 2}

\cas{1)}
\explique{$x^{2}+x-12=0$ : $a=1$, $b=1$ et $c=-12$.}{}
\begin{calculs}
\Delta &= 1^{2}-4\times 1\times(-12) &&\justif{$\Delta=b^{2}-4ac$}\\
\Delta &= 1+48 &&\justif{$-4\times(-12)=+48$}\\
\Delta &= 49=7^{2} &&\justif{$\Delta>0$ : deux solutions}\\
x_{1} &= \dfrac{-1+7}{2}=3 &&\justif{$x_{1}=\frac{-b+\sqrt{\Delta}}{2a}$}\\
x_{2} &= \dfrac{-1-7}{2}=-4 &&\justif{$x_{2}=\frac{-b-\sqrt{\Delta}}{2a}$}
\end{calculs}
\explique{$S=\res{\{-4\,;3\}}$.}{contrôle : $3^{2}+3-12=0$}
\cas{2)}
\explique{$x^{2}+3x+5=0$ : $a=1$, $b=3$ et $c=5$.}{}
\begin{calculs}
\Delta &= 3^{2}-4\times 1\times 5 &&\\
\Delta &= 9-20=-11 &&\justif{$\Delta<0$ : aucune solution réelle}
\end{calculs}
\explique{$S=\res{\varnothing}$.}{on écrit l'ensemble vide, pas « $S=0$ »}
\cas{3)}
\explique{$-x^{2}+6x-9=0$ : $a=-1$, $b=6$ et $c=-9$.}{}
\begin{calculs}
\Delta &= 6^{2}-4\times(-1)\times(-9) &&\justif{$(-1)\times(-9)=9$ : on retire $4\times 9$}\\
\Delta &= 36-36=0 &&\justif{$\Delta=0$ : une solution (racine double)}\\
x_{0} &= -\dfrac{6}{2\times(-1)}=3 &&\justif{$x_{0}=-\frac{b}{2a}$}
\end{calculs}
\explique{$S=\res{\{3\}}$.}{}
\rem{Le trinôme vaut $-\left(x^{2}-6x+9\right)=-(x-3)^{2}$ : on pouvait factoriser
directement.}
\cas{4)}
\explique{L'inconnue s'appelle $t$ : rien ne change. $a=1$, $b=-4$ et $c=-21$.}{}
\begin{calculs}
\Delta &= (-4)^{2}-4\times 1\times(-21) &&\justif{$(-4)^{2}=16$, avec les parenthèses}\\
\Delta &= 16+84 &&\\
\Delta &= 100=10^{2} &&\justif{$\Delta>0$ : deux solutions}\\
t_{1} &= \dfrac{4+10}{2}=7 &&\justif{$-b=4$}\\
t_{2} &= \dfrac{4-10}{2}=-3 &&
\end{calculs}
\explique{$S=\res{\{-3\,;7\}}$.}{}
\cas{5)}
\explique{On range d'abord les termes dans l'ordre $ay^{2}+by+c$ :
$2-y-3y^{2}=0 \iff -3y^{2}-y+2=0$. Donc $a=-3$, $b=-1$ et $c=2$.}{$a$ est le
coefficient de $y^{2}$, quelle que soit sa place}
\begin{calculs}
\Delta &= (-1)^{2}-4\times(-3)\times 2 &&\\
\Delta &= 1+24 &&\\
\Delta &= 25=5^{2} &&\justif{$\Delta>0$ : deux solutions}\\
y_{1} &= \dfrac{1+5}{-6}=-1 &&\justif{$-b=1$ et $2a=-6$}\\
y_{2} &= \dfrac{1-5}{-6}=\dfrac{2}{3} &&\justif{$\frac{-4}{-6}=\frac{2}{3}$}
\end{calculs}
\explique{$S=\res{\left\{-1\,;\dfrac{2}{3}\right\}}$.}{}
\rem{Erreur fréquente : lire $a=2$ parce que $2$ est écrit en premier.}
\cas{6)}
\explique{$-2x^{2}+5x+1=0$ : $a=-2$, $b=5$ et $c=1$.}{}
\begin{calculs}
\Delta &= 5^{2}-4\times(-2)\times 1 &&\\
\Delta &= 25+8=33 &&\justif{$\Delta>0$, mais $33$ n'est pas un carré : on garde $\sqrt{33}$}\\
x_{1} &= \dfrac{-5+\sqrt{33}}{-4}=\dfrac{5-\sqrt{33}}{4} &&\justif{diviser par $-4$ : on change tous les signes}\\
x_{2} &= \dfrac{-5-\sqrt{33}}{-4}=\dfrac{5+\sqrt{33}}{4} &&
\end{calculs}
\explique{$S=\res{\left\{\dfrac{5-\sqrt{33}}{4}\,;\dfrac{5+\sqrt{33}}{4}\right\}}$.}{valeurs
exactes ; à la calculatrice, $x_{1}\approx -0{,}19$ et $x_{2}\approx 2{,}69$}
\end{exo}

\begin{exo}[Des racines carrées dans les coefficients]
\cas{1)}
\explique{$x^{2}-2\sqrt{5}\,x-4=0$ : $a=1$, $b=-2\sqrt{5}$ et $c=-4$.}{la méthode ne
change pas avec des racines carrées}
\begin{calculs}
\Delta &= \left(-2\sqrt{5}\right)^{2}-4\times 1\times(-4) &&\justif{$\Delta=b^{2}-4ac$}\\
\Delta &= 20+16 &&\justif{$\left(-2\sqrt{5}\right)^{2}=4\times 5=20$}\\
\Delta &= 36=6^{2} &&\justif{$\Delta>0$ : deux solutions}\\
x_{1} &= \dfrac{2\sqrt{5}+6}{2}=\sqrt{5}+3 &&\justif{on divise chaque terme par $2$}\\
x_{2} &= \dfrac{2\sqrt{5}-6}{2}=\sqrt{5}-3 &&
\end{calculs}
\explique{$S=\res{\left\{\sqrt{5}-3\,;\sqrt{5}+3\right\}}$.}{$\sqrt{5}-3<0<\sqrt{5}+3$}
\cas{2)}
\explique{$\sqrt{3}\,t^{2}-4t+\sqrt{3}=0$ : $a=\sqrt{3}$, $b=-4$ et
$c=\sqrt{3}$.}{}
\begin{calculs}
\Delta &= (-4)^{2}-4\times\sqrt{3}\times\sqrt{3} &&\\
\Delta &= 16-12 &&\justif{$\sqrt{3}\times\sqrt{3}=3$}\\
\Delta &= 4=2^{2} &&\justif{$\Delta>0$ : deux solutions}\\
t_{1} &= \dfrac{4+2}{2\sqrt{3}}=\dfrac{3}{\sqrt{3}} &&\justif{on simplifie par $2$}\\
t_{1} &= \dfrac{\sqrt{3}\times\barre{\sqrt{3}}}{\barre{\sqrt{3}}}=\sqrt{3} &&\justif{$3=\sqrt{3}\times\sqrt{3}$}\\
t_{2} &= \dfrac{4-2}{2\sqrt{3}}=\dfrac{1}{\sqrt{3}} &&\\
t_{2} &= \dfrac{1\fois{\sqrt{3}}}{\sqrt{3}\fois{\sqrt{3}}}=\dfrac{\sqrt{3}}{3} &&\justif{sans racine au dénominateur}
\end{calculs}
\explique{$S=\res{\left\{\dfrac{\sqrt{3}}{3}\,;\sqrt{3}\right\}}$.}{$\frac{1}{\sqrt{3}}$
est aussi une réponse juste}
\cas{3)}
\explique{$x^{2}-(1+\sqrt{2})x+\sqrt{2}=0$ : $a=1$, $b=-(1+\sqrt{2})$ et
$c=\sqrt{2}$.}{$b$ est tout le facteur devant $x$, avec son signe}
\begin{calculs}
\Delta &= (1+\sqrt{2})^{2}-4\sqrt{2} &&\justif{$b^{2}=(1+\sqrt{2})^{2}$ : le signe disparaît au carré}\\
\Delta &= 1+2\sqrt{2}+2-4\sqrt{2} &&\justif{identité remarquable}\\
\Delta &= 3-2\sqrt{2} &&\justif{$\Delta>0$ car $2\sqrt{2}\approx 2{,}83$}
\end{calculs}
\explique{Pour calculer $\sqrt{\Delta}$, on cherche à écrire $3-2\sqrt{2}$ comme un carré
: le « $-2\sqrt{2}$ » fait penser au double produit de $(\sqrt{2}-1)^{2}$.}{c'est
l'étape qui n'a rien d'automatique}
\begin{calculs}
(\sqrt{2}-1)^{2} &= 2-2\sqrt{2}+1 &&\justif{$(u-v)^{2}$ avec $u=\sqrt{2}$ et $v=1$}\\
(\sqrt{2}-1)^{2} &= 3-2\sqrt{2} &&\justif{c'est bien $\Delta$ ; et $\sqrt{2}-1>0$, donc $\sqrt{\Delta}=\sqrt{2}-1$}\\
x_{1} &= \dfrac{1+\sqrt{2}+(\sqrt{2}-1)}{2}=\dfrac{2\sqrt{2}}{2} &&\justif{$-b=1+\sqrt{2}$}\\
x_{1} &= \sqrt{2} &&\\
x_{2} &= \dfrac{1+\sqrt{2}-(\sqrt{2}-1)}{2}=\dfrac{2}{2} &&\justif{les $\sqrt{2}$ s'éliminent}\\
x_{2} &= 1 &&
\end{calculs}
\explique{$S=\res{\{1\,;\sqrt{2}\}}$.}{contrôle :
$1-(1+\sqrt{2})+\sqrt{2}=0$}
\cas{4)}
\explique{On range : $4x-x^{2}-5=0 \iff -x^{2}+4x-5=0$, donc $a=-1$, $b=4$ et
$c=-5$.}{même piège qu'à l'exercice 8, question 5)}
\begin{calculs}
\Delta &= 4^{2}-4\times(-1)\times(-5) &&\\
\Delta &= 16-20=-4 &&\justif{$\Delta<0$ : aucune solution}
\end{calculs}
\explique{$S=\res{\varnothing}$.}{}
\cas{5)}
\explique{L'équation est de degré $3$ : pas de discriminant. Mais $x$ est en facteur dans
chaque terme.}{on se ramène à un produit nul}
\begin{calculs}
x^{3}+2x^{2}-15x=0 &\iff x\left(x^{2}+2x-15\right)=0 &&\justif{facteur commun $x$}\\
&\iff x=0 \text{ ou } x^{2}+2x-15=0 &&\justif{produit nul}
\end{calculs}
\begin{calculs}
\Delta &= 2^{2}-4\times 1\times(-15) &&\justif{pour $x^{2}+2x-15=0$}\\
\Delta &= 64=8^{2} &&\\
x_{1} &= \dfrac{-2+8}{2}=3 &&\\
x_{2} &= \dfrac{-2-8}{2}=-5 &&
\end{calculs}
\explique{$S=\res{\{-5\,;0\,;3\}}$.}{trois solutions : on n'oublie pas $0$}
\cas{6)}
\explique{Pas besoin de développer : $(3x+2)^{2}\ge 0$ pour tout réel $x$, donc
$(3x+2)^{2}+1\ge 1$. Le membre de gauche ne vaut jamais $0$.}{un carré est positif}
\explique{$S=\res{\varnothing}$.}{}
\rem{En développant, on obtient $9x^{2}+12x+5=0$ et $\Delta=144-180=-36<0$ : même
conclusion, en plus long. En 5) et 6), regarder la forme de l'équation avant de
calculer.}
\end{exo}

\begin{exo}[Une parabole et une droite]
\cas{1)}
\begin{minipage}[c]{0.56\linewidth}\raggedright
\explique{Parabole $y=x^{2}$ : on place quelques points, $(0\,;0)$, $(1\,;1)$,
$(-1\,;1)$, $(2\,;4)$, $(-2\,;4)$ et $(3\,;9)$.}{courbe symétrique par rapport à l'axe
des ordonnées}
\explique{Droite $y=2x+3$ : deux points suffisent, $(0\,;3)$ et $(3\,;9)$.}{ordonnée à
l'origine $3$, coefficient directeur $2$}
\end{minipage}\hfill
\begin{minipage}[c]{0.4\linewidth}\centering
\begin{tikzpicture}
\begin{axis}[repere,x=0.45cm,y=0.45cm,
    xmin=-2.5,xmax=4,ymin=-0.8,ymax=10,
    xtick={-2,...,3},ytick={1,...,9}]
  \addplot[courbe,domain=-2.4:3.15]{x^2};
  \addplot[coulCourbeB,line width=1pt,domain=-1.9:3.4]{2*x+3};
  \addplot[croix] coordinates {(-1,1)(3,9)};
\end{axis}
\end{tikzpicture}
\end{minipage}

\cas{2)}
\explique{Un point $M(x\,;y)$ est commun aux deux courbes quand $y=x^{2}$ \emph{et}
$y=2x+3$, c'est-à-dire quand $x^{2}=2x+3$.}{même ordonnée pour la même abscisse}
\begin{calculs}
x^{2}=2x+3 &\iff x^{2}-2x-3=0 &&\justif{on passe tout dans le membre de gauche}
\end{calculs}
\explique{Les solutions de $x^{2}-2x-3=0$ sont donc les abscisses des points communs.
Sur la figure, il y a \res{deux} points communs : l'équation a deux solutions, qu'on lit
\res{$-1$ et $3$}.}{une lecture graphique se vérifie ensuite par le calcul}
\cas{3)}
\explique{$a=1$, $b=-2$ et $c=-3$.}{}
\begin{calculs}
\Delta &= (-2)^{2}-4\times 1\times(-3) &&\\
\Delta &= 4+12=16 &&\justif{$\Delta>0$ : deux solutions, comme sur la figure}\\
x_{1} &= \dfrac{2+4}{2}=3 &&\\
x_{2} &= \dfrac{2-4}{2}=-1 &&
\end{calculs}
\explique{$S=\res{\{-1\,;3\}}$ : la lecture est confirmée.}{contrôle :
$(-1)^{2}=1$ et $2\times(-1)+3=1$ ; $3^{2}=9$ et $2\times 3+3=9$}
\end{exo}

\begin{exo}[Un programme Python]
\cas{1)}
\explique{\texttt{d = b*b - 4*a*c} calcule $b^{2}-4ac$ : \texttt{d} est le
\res{discriminant $\Delta$}.}{\texttt{b*b} est $b^{2}$}
\explique{Si \texttt{d < 0}, la fonction renvoie la liste vide \texttt{[]} : aucune
solution. Si \texttt{d == 0}, elle renvoie la seule solution $-\dfrac{b}{2a}$. Sinon,
elle calcule $r=\sqrt{\Delta}$ et renvoie les deux solutions
$\dfrac{-b-\sqrt{\Delta}}{2a}$ et $\dfrac{-b+\sqrt{\Delta}}{2a}$.}{ce sont exactement
les trois cas du théorème 2}
\cas{2)}
\explique{On entre $a$, $b$ et $c$ de chaque équation, une fois ordonnée.}{en 5), on
tape \texttt{racines(-3, -1, 2)}, pas \texttt{racines(2, -1, -3)}}
\begin{calculs}
&\texttt{racines(1, 1, -12)} \text{ renvoie } \texttt{[-4.0, 3.0]} &&\justif{1)}\\
&\texttt{racines(1, 3, 5)} \text{ renvoie } \texttt{[]} &&\justif{2) : aucune solution}\\
&\texttt{racines(-1, 6, -9)} \text{ renvoie } \texttt{[3.0]} &&\justif{3) : une solution}\\
&\texttt{racines(1, -4, -21)} \text{ renvoie } \texttt{[-3.0, 7.0]} &&\justif{4)}\\
&\texttt{racines(-3, -1, 2)} \text{ renvoie } \texttt{[0.666\dots, -1.0]} &&\justif{5) : $0{,}666\ldots\approx\frac{2}{3}$}\\
&\texttt{racines(-2, 5, 1)} \text{ renvoie } \texttt{[2.686\dots, -0.186\dots]} &&\justif{6) : valeurs approchées}
\end{calculs}
\explique{On retrouve les résultats obtenus à la main ; en 6),
$\dfrac{5+\sqrt{33}}{4}\approx 2{,}686$ et $\dfrac{5-\sqrt{33}}{4}\approx -0{,}186$.}{le
programme donne des valeurs approchées, la main des valeurs exactes}
\rem{Quand $a<0$, le programme range les solutions dans l'ordre décroissant : diviser
par $2a<0$ inverse l'ordre.}
\cas{3)}
\explique{Avec $a=0$, $b=3$ et $c=-12$ : \texttt{d} vaut $9-0=9>0$, et le programme
calcule $\dfrac{-3-3}{2\times 0}$. Python s'arrête sur une \res{erreur de division par
zéro} (\texttt{ZeroDivisionError}).}{on ne divise jamais par $0$}
\explique{Avec $a=0$, l'équation $3x-12=0$ n'est pas du second degré : les formules du
discriminant ne s'appliquent pas. Elle se résout directement :
$3x-12=0 \iff x=\res{4}$.}{un trinôme exige $a\ne 0$ (définition 1)}
\cas{4)}
\explique{\texttt{racines(1, 6, 9)} renvoie \texttt{[-3.0]} : $\Delta=36-36=0$ est
calculé exactement.}{nombres entiers : calcul exact}
\explique{$0{,}2(x+3)^{2}$ a aussi la racine double $-3$ : on attendait
\texttt{[-3.0]}. Le programme renvoie pourtant \res{\texttt{[]}}.}{$\Delta=1{,}2^{2}-4\times 0{,}2\times 1{,}8=0$ en théorie}
\explique{L'ordinateur enregistre $0{,}2$ ; $1{,}2$ et $1{,}8$ en valeur
\emph{approchée} (en binaire) : \texttt{d} vaut environ $-2{,}2\times 10^{-16}$ au lieu
de $0$. Le test \texttt{d < 0} est alors vrai, et le programme conclut à tort qu'il n'y a
pas de solution.}{une erreur minuscule suffit à changer le signe d'un nombre nul}
\rem{Un test d'égalité \texttt{d == 0} sur des nombres à virgule calculés est fragile ;
on teste plutôt si \texttt{d} est très proche de $0$, par exemple
\texttt{abs(d) < 1e-12}. Le résultat exact peut dépendre de la machine.}
\end{exo}

\partie{Factorisation du trinôme}

\begin{exo}[Factoriser un trinôme]
\explique{On cherche les racines avec le discriminant, puis on écrit
$a(x-x_{1})(x-x_{2})$, sans oublier $a$ devant.}{méthode 4 et théorème 3}

\cas{1)}
\explique{$f(x)=x^{2}+3x-10$ : $a=1$, $b=3$ et $c=-10$.}{}
\begin{calculs}
\Delta &= 3^{2}-4\times 1\times(-10) &&\\
\Delta &= 9+40 &&\\
\Delta &= 49=7^{2} &&\justif{$\Delta>0$ : deux racines, $f$ se factorise}\\
x_{1} &= \dfrac{-3+7}{2}=2 &&\\
x_{2} &= \dfrac{-3-7}{2}=-5 &&\\
f(x) &= (x-2)\bigl(x-(-5)\bigr) &&\justif{$a(x-x_{1})(x-x_{2})$ avec $a=1$}\\
f(x) &= \res{(x-2)(x+5)} &&\justif{$x-(-5)=x+5$}
\end{calculs}
\explique{Contrôle : le terme constant de $(x-2)(x+5)$ est $(-2)\times 5=-10$, c'est bien
$c$.}{$a\,x_{1}x_{2}$ doit valoir $c$}

\cas{2)}
\explique{$g(x)=3x^{2}-10x+3$ : $a=3$, $b=-10$ et $c=3$.}{}
\begin{calculs}
\Delta &= (-10)^{2}-4\times 3\times 3 &&\\
\Delta &= 100-36 &&\\
\Delta &= 64=8^{2} &&\justif{$\Delta>0$ : deux racines}\\
x_{1} &= \dfrac{10+8}{6}=3 &&\justif{$-b=10$ et $2a=6$}\\
x_{2} &= \dfrac{10-8}{6}=\dfrac{1}{3} &&\\
g(x) &= 3(x-3)\left(x-\dfrac{1}{3}\right) &&\justif{on n'oublie pas $a=3$ devant}\\
g(x) &= \res{(x-3)(3x-1)} &&\justif{le $3$ entre dans la seconde parenthèse : $3\left(x-\frac{1}{3}\right)=3x-1$}
\end{calculs}
\explique{Contrôle : terme constant $(-3)\times(-1)=3=c$.}{}

\cas{3)}
\explique{$h(x)=-2x^{2}+3x+5$ : $a=-2$, $b=3$ et $c=5$.}{}
\begin{calculs}
\Delta &= 3^{2}-4\times(-2)\times 5 &&\\
\Delta &= 9+40 &&\\
\Delta &= 49=7^{2} &&\justif{$\Delta>0$ : deux racines}\\
x_{1} &= \dfrac{-3+7}{-4}=-1 &&\justif{$2a=-4$}\\
x_{2} &= \dfrac{-3-7}{-4}=\dfrac{5}{2} &&\justif{$\frac{-10}{-4}=\frac{5}{2}$}\\
h(x) &= -2(x+1)\left(x-\dfrac{5}{2}\right) &&\justif{$a=-2$ devant ; $x-(-1)=x+1$}\\
h(x) &= \res{-(x+1)(2x-5)} &&\justif{$-2=-1\times 2$ : le $2$ entre dans la seconde parenthèse}
\end{calculs}
\explique{Contrôle : terme constant $-2\times(-1)\times\dfrac{5}{2}=5=c$.}{$a\,x_{1}x_{2}=c$}

\cas{4)}
\explique{Les coefficients sont des fractions : on met d'abord $a=-\dfrac{1}{3}$ en
facteur. La parenthèse a alors des coefficients entiers, et les mêmes racines.}{comme
pour la forme canonique quand $a\ne 1$}
\begin{calculs}
k(x) &= -\dfrac{1}{3}x^{2}+\dfrac{1}{3}x+2 &&\justif{on recopie}\\
k(x) &= -\dfrac{1}{3}\left(x^{2}-x-6\right) &&\justif{$\frac{1}{3}x=-\frac{1}{3}\times(-x)$ et $2=-\frac{1}{3}\times(-6)$}
\end{calculs}
\begin{calculs}
\Delta &= (-1)^{2}-4\times 1\times(-6) &&\justif{pour $x^{2}-x-6$}\\
\Delta &= 25=5^{2} &&\\
x_{1} &= \dfrac{1+5}{2}=3 &&\\
x_{2} &= \dfrac{1-5}{2}=-2 &&\\
k(x) &= \res{-\dfrac{1}{3}(x-3)(x+2)} &&\justif{$x^{2}-x-6=(x-3)(x+2)$}
\end{calculs}
\explique{Contrôle : terme constant $-\dfrac{1}{3}\times(-3)\times 2=2=c$.}{}
\rem{Erreur fréquente : écrire $(x-x_{1})(x-x_{2})$ sans le facteur $a$. Le contrôle
par le terme constant la détecte : sans $a$, en 3), on trouverait $-\frac{5}{2}$ au
lieu de $5$.}
\end{exo}

\begin{exo}[Factoriser, ou prouver que c'est impossible]
\explique{Si $\Delta<0$, le trinôme n'a pas de racine et ne se factorise pas : c'est
la preuve demandée. Une identité remarquable se repère avant tout calcul.}{théorème 3}

\cas{A}
\explique{$A=x^{2}+2x+2$ : $a=1$, $b=2$ et $c=2$.}{}
\begin{calculs}
\Delta &= 2^{2}-4\times 1\times 2 &&\\
\Delta &= 4-8=-4 &&\justif{$\Delta<0$ : aucune racine}
\end{calculs}
\explique{Le trinôme $A$ \res{ne se factorise pas}.}{on cite le signe de $\Delta$ :
c'est la démonstration}
\cas{B}
\begin{calculs}
B &= x^{2}+8x+16 &&\justif{$x^{2}+2\times x\times 4+4^{2}$}\\
B &= \res{(x+4)^{2}} &&\justif{$(u+v)^{2}$ : pas de $\Delta$}
\end{calculs}
\explique{Avec le discriminant : $\Delta=64-64=0$ et $x_{0}=-4$, donc
$B=(x+4)^{2}$.}{même résultat, plus long}
\cas{C}
\explique{$C=3x^{2}-5x+2$ : $a=3$, $b=-5$ et $c=2$.}{}
\begin{calculs}
\Delta &= (-5)^{2}-4\times 3\times 2 &&\\
\Delta &= 25-24=1 &&\justif{$\Delta>0$ : deux racines}\\
x_{1} &= \dfrac{5+1}{6}=1 &&\\
x_{2} &= \dfrac{5-1}{6}=\dfrac{2}{3} &&\\
C &= 3(x-1)\left(x-\dfrac{2}{3}\right) &&\justif{$a=3$ devant}\\
C &= \res{(x-1)(3x-2)} &&\justif{le $3$ entre dans la seconde parenthèse}
\end{calculs}
\cas{D}
\explique{$D=-x^{2}+4x-7$ : $a=-1$, $b=4$ et $c=-7$.}{}
\begin{calculs}
\Delta &= 4^{2}-4\times(-1)\times(-7) &&\justif{$(-1)\times(-7)=7$}\\
\Delta &= 16-28=-12 &&\justif{$\Delta<0$}
\end{calculs}
\explique{Le trinôme $D$ \res{ne se factorise pas}.}{}
\cas{E}
\begin{calculs}
E &= 9x^{2}+12x+4 &&\justif{$(3x)^{2}+2\times 3x\times 2+2^{2}$}\\
E &= \res{(3x+2)^{2}} &&\justif{$(u+v)^{2}$ avec $u=3x$ et $v=2$}
\end{calculs}
\cas{F}
\explique{$F=2x^{2}+3x+1$ : $a=2$, $b=3$ et $c=1$.}{}
\begin{calculs}
\Delta &= 3^{2}-4\times 2\times 1 &&\\
\Delta &= 9-8=1 &&\\
x_{1} &= \dfrac{-3+1}{4}=-\dfrac{1}{2} &&\\
x_{2} &= \dfrac{-3-1}{4}=-1 &&\\
F &= 2\left(x+\dfrac{1}{2}\right)(x+1) &&\justif{$a=2$ devant}\\
F &= \res{(2x+1)(x+1)} &&\justif{le $2$ entre dans la première parenthèse}
\end{calculs}
\rem{« Démontrer qu'il ne se factorise pas » : un calcul de $\Delta<0$ et la phrase de
conclusion suffisent ; ne pas se contenter de « je n'y arrive pas ».}
\end{exo}

\begin{exo}[Simplifier un quotient]
\explique{On écrit d'abord l'ensemble de définition (dénominateur non nul), puis on
factorise pour faire apparaître un facteur commun en haut et en bas.}{on ne simplifie
que des \emph{facteurs}}

\cas{1)}
\explique{$x+3=0 \iff x=-3$, donc $D=\res{\R\setminus\{-3\}}$.}{valeur interdite}
\explique{Numérateur $x^{2}+x-6$ : $a=1$, $b=1$ et $c=-6$.}{}
\begin{calculs}
\Delta &= 1^{2}-4\times 1\times(-6) &&\\
\Delta &= 25=5^{2} &&\\
x_{1} &= \dfrac{-1+5}{2}=2 &&\\
x_{2} &= \dfrac{-1-5}{2}=-3 &&\justif{la valeur interdite est aussi racine du numérateur}\\
A(x) &= \dfrac{(x-2)\barre{(x+3)}}{\barre{x+3}} &&\justif{$x^{2}+x-6=(x-2)(x+3)$}\\
A(x) &= \res{x-2} &&\justif{pour $x\ne -3$}
\end{calculs}
\cas{2)}
\explique{$x^{2}-1=(x-1)(x+1)$ s'annule en $1$ et en $-1$, donc
$D=\res{\R\setminus\{-1\,;1\}}$.}{$u^{2}-v^{2}$}
\explique{Numérateur $3x^{2}+x-2$ : $a=3$, $b=1$ et $c=-2$.}{}
\begin{calculs}
\Delta &= 1^{2}-4\times 3\times(-2) &&\\
\Delta &= 25=5^{2} &&\\
x_{1} &= \dfrac{-1+5}{6}=\dfrac{2}{3} &&\\
x_{2} &= \dfrac{-1-5}{6}=-1 &&\\
3x^{2}+x-2 &= 3\left(x-\dfrac{2}{3}\right)(x+1) &&\justif{$a=3$ devant}\\
3x^{2}+x-2 &= (3x-2)(x+1) &&\justif{le $3$ entre dans la première parenthèse}\\
B(x) &= \dfrac{(3x-2)\barre{(x+1)}}{(x-1)\barre{(x+1)}} &&\justif{facteur commun $x+1$}\\
B(x) &= \res{\dfrac{3x-2}{x-1}} &&\justif{pour $x\in D$}
\end{calculs}
\rem{L'égalité simplifiée n'est vraie que sur $D$ : en 1), $A(-3)$ n'existe pas, alors
que $x-2$ vaut $-5$ en $-3$. Erreur fréquente : « simplifier » des termes, comme
$\frac{x^{2}+x-6}{x+3}$ en barrant les $x$.}
\end{exo}

\partie{Somme et produit des racines}

\begin{exo}[Une racine connue (1)]
\cas{1)}
\begin{calculs}
(-2)^{2}+5\times(-2)+6 &= 4-10+6 &&\justif{on remplace $x$ par $-2$}\\
(-2)^{2}+5\times(-2)+6 &= 0 &&
\end{calculs}
\explique{\res{$-2$ est solution} de l'équation.}{une valeur est solution si elle rend
l'égalité vraie}
\cas{2)}
\explique{$a=1$, $b=5$ et $c=6$.}{théorème 4}
\begin{calculs}
S &= -\dfrac{b}{a}=\res{-5} &&\justif{attention au signe « $-$ » de la formule}\\
P &= \dfrac{c}{a}=\res{6} &&
\end{calculs}
\cas{3)}
\begin{calculs}
x_{2} &= S-x_{1} &&\justif{$x_{1}+x_{2}=S$}\\
x_{2} &= -5-(-2) &&\justif{on retranche $-2$ : parenthèses}\\
x_{2} &= \res{-3} &&
\end{calculs}
\explique{Contrôle : $(-2)\times(-3)=6$, c'est bien $P$. Les solutions sont $-3$ et
$-2$.}{le produit confirme le calcul fait avec la somme}
\end{exo}

\begin{exo}[Une racine connue (2)]
\cas{1)}
\begin{calculs}
3^{2}-7\times 3+12 &= 9-21+12 &&\justif{on remplace $x$ par $3$}\\
3^{2}-7\times 3+12 &= 0 &&
\end{calculs}
\explique{\res{$3$ est solution}.}{}
\cas{2)}
\explique{$a=1$, $b=-7$ et $c=12$ : $S=-\dfrac{-7}{1}=\res{7}$ et
$P=\dfrac{12}{1}=\res{12}$.}{}
\cas{3)}
\begin{calculs}
x_{2} &= S-x_{1} &&\\
x_{2} &= 7-3=\res{4} &&
\end{calculs}
\explique{Contrôle : $3\times 4=12=P$. Les solutions sont $3$ et $4$.}{}
\end{exo}

\begin{exo}[Racine évidente, puis somme]
\explique{On essaie $1$, puis $-1$, puis $2$\dots{} de tête ; une fois $x_{1}$ trouvée,
$x_{2}=S-x_{1}$ avec $S=-\dfrac{b}{a}$.}{méthode 5}

\cas{1)}
\begin{calculs}
1^{2}-6\times 1+5 &= 0 &&\justif{$x_{1}=1$ est racine évidente}\\
S &= -\dfrac{-6}{1}=6 &&\justif{$S=-\frac{b}{a}$}\\
x_{2} &= 6-1=\res{5} &&\justif{$x_{2}=S-x_{1}$}
\end{calculs}
\explique{Contrôle : $1\times 5=5=\dfrac{c}{a}$. Les racines sont $1$ et $5$.}{}
\cas{2)}
\begin{calculs}
4\times 1^{2}+1-5 &= 0 &&\justif{$x_{1}=1$}\\
S &= -\dfrac{1}{4} &&\\
x_{2} &= -\dfrac{1}{4}-\dfrac{1\fois{4}}{1\fois{4}} &&\justif{$x_{2}=S-1$ ; même dénominateur}\\
x_{2} &= \res{-\dfrac{5}{4}} &&
\end{calculs}
\explique{Contrôle : $1\times\left(-\dfrac{5}{4}\right)=-\dfrac{5}{4}=\dfrac{c}{a}$.}{}
\cas{3)}
\begin{calculs}
7\times 1^{2}-9\times 1+2 &= 0 &&\justif{$x_{1}=1$}\\
S &= -\dfrac{-9}{7}=\dfrac{9}{7} &&\\
x_{2} &= \dfrac{9}{7}-\dfrac{1\fois{7}}{1\fois{7}} &&\\
x_{2} &= \res{\dfrac{2}{7}} &&
\end{calculs}
\explique{Contrôle : $1\times\dfrac{2}{7}=\dfrac{2}{7}=\dfrac{c}{a}$.}{}
\cas{4)}
\begin{calculs}
3\times 1^{2}+8\times 1+5 &= 16 &&\justif{$1$ ne convient pas}\\
3\times(-1)^{2}+8\times(-1)+5 &= 0 &&\justif{$x_{1}=-1$}\\
S &= -\dfrac{8}{3} &&\\
x_{2} &= -\dfrac{8}{3}-(-1) &&\justif{on retranche $-1$}\\
x_{2} &= -\dfrac{8}{3}+\dfrac{1\fois{3}}{1\fois{3}} &&\\
x_{2} &= \res{-\dfrac{5}{3}} &&
\end{calculs}
\explique{Contrôle : $(-1)\times\left(-\dfrac{5}{3}\right)=\dfrac{5}{3}=\dfrac{c}{a}$.}{}
\rem{En remplaçant $x$ par $1$, on obtient $a+b+c$ : si $a+b+c=0$, alors $1$ est racine. En
remplaçant par $-1$, on obtient $a-b+c$ : si $a-b+c=0$, alors $-1$ est racine.}
\end{exo}

\begin{exo}[Sans discriminant]
\cas{1)}
\begin{calculs}
1-10+9 &= 0 &&\justif{$x_{1}=1$ est racine évidente}\\
S &= 10 &&\justif{$S=-\frac{b}{a}$}\\
x_{2} &= 10-1=9 &&
\end{calculs}
\explique{Contrôle : $1\times 9=9=\dfrac{c}{a}$. Les solutions sont \res{$1$ et
$9$}.}{}
\cas{2)}
\begin{calculs}
-2+3+5 &= 6 &&\justif{$1$ ne convient pas}\\
-2\times(-1)^{2}+3\times(-1)+5 &= 0 &&\justif{$x_{1}=-1$}\\
S &= -\dfrac{3}{-2}=\dfrac{3}{2} &&\\
x_{2} &= \dfrac{3}{2}+\dfrac{1\fois{2}}{1\fois{2}} &&\justif{$x_{2}=S-(-1)=S+1$}\\
x_{2} &= \dfrac{5}{2} &&
\end{calculs}
\explique{Contrôle : $(-1)\times\dfrac{5}{2}=-\dfrac{5}{2}=\dfrac{5}{-2}$. Les solutions
sont \res{$-1$ et $\dfrac{5}{2}$}.}{$\frac{c}{a}=\frac{5}{-2}$}
\cas{3)}
\explique{$1$ ne convient pas ($1+4-12=-7$), $-1$ non plus ($1-4-12=-15$) ; $2$
convient : $4+8-12=0$.}{on essaie dans l'ordre $1$ ; $-1$ ; $2$\dots}
\begin{calculs}
S &= -4 &&\\
x_{2} &= -4-2=-6 &&
\end{calculs}
\explique{Contrôle : $2\times(-6)=-12=\dfrac{c}{a}$. Les solutions sont \res{$-6$ et
$2$}.}{}
\cas{4)}
\begin{calculs}
1-7+6 &= 0 &&\justif{$x_{1}=-1$ : $(-1)^{2}+7\times(-1)+6$}\\
S &= -7 &&\\
x_{2} &= -7-(-1)=-6 &&
\end{calculs}
\explique{Contrôle : $(-1)\times(-6)=6=\dfrac{c}{a}$. Les solutions sont \res{$-6$ et
$-1$}.}{}
\cas{5)}
\explique{$1$ ; $-1$ ; $2$ et $-2$ ne conviennent pas ; $3$ convient :
$9-27+18=0$.}{on continue les essais}
\begin{calculs}
S &= 9 &&\\
x_{2} &= 9-3=6 &&
\end{calculs}
\explique{Contrôle : $3\times 6=18=\dfrac{c}{a}$. Les solutions sont \res{$3$ et
$6$}.}{deux nombres de somme $9$ et de produit $18$}
\cas{6)}
\explique{Le $\sqrt{2}$ de l'équation invite à essayer $\sqrt{2}$ :
$\left(\sqrt{2}\right)^{2}+\sqrt{2}\times\sqrt{2}-4=2+2-4=0$.}{$x_{1}=\sqrt{2}$}
\begin{calculs}
S &= -\sqrt{2} &&\justif{$a=1$ et $b=\sqrt{2}$}\\
x_{2} &= -\sqrt{2}-\sqrt{2} &&\\
x_{2} &= -2\sqrt{2} &&
\end{calculs}
\explique{Contrôle : $\sqrt{2}\times\left(-2\sqrt{2}\right)=-4=\dfrac{c}{a}$. Les
solutions sont \res{$-2\sqrt{2}$ et $\sqrt{2}$}.}{}
\rem{Quand aucune valeur simple ne marche après quelques essais, on revient au
discriminant : la méthode 5 est un raccourci, pas une obligation.}
\end{exo}

\begin{exo}[Calculer sans résoudre]
\cas{1)}
\explique{$2x^{2}-9x+4=0$ : $a=2$, $b=-9$ et $c=4$.}{}
\begin{calculs}
\Delta &= (-9)^{2}-4\times 2\times 4 &&\\
\Delta &= 81-32=49 &&\justif{$\Delta>0$ : les deux racines existent}\\
S &= -\dfrac{-9}{2}=\res{\dfrac{9}{2}} &&\justif{$S=-\frac{b}{a}$}\\
P &= \dfrac{4}{2}=\res{2} &&\justif{$P=\frac{c}{a}$}
\end{calculs}
\cas{2)}
\explique{$P=2\ne 0$, donc aucune racine n'est nulle : on peut diviser par $x_{1}$ et
$x_{2}$.}{}
\begin{calculs}
\dfrac{1}{x_{1}}+\dfrac{1}{x_{2}} &= \dfrac{1\fois{x_{2}}}{x_{1}\fois{x_{2}}}+\dfrac{1\fois{x_{1}}}{x_{2}\fois{x_{1}}} &&\justif{même dénominateur $x_{1}x_{2}$}\\
\dfrac{1}{x_{1}}+\dfrac{1}{x_{2}} &= \dfrac{x_{1}+x_{2}}{x_{1}x_{2}} &&\justif{on reconnaît $S$ en haut et $P$ en bas}\\
\dfrac{1}{x_{1}}+\dfrac{1}{x_{2}} &= \dfrac{9}{2}\times\dfrac{1}{2} &&\justif{$\frac{S}{P}$ ; diviser par $2$, c'est multiplier par $\frac{1}{2}$}\\
\dfrac{1}{x_{1}}+\dfrac{1}{x_{2}} &= \res{\dfrac{9}{4}} &&
\end{calculs}
\cas{3)}
\begin{calculs}
S^{2} &= (x_{1}+x_{2})^{2} &&\\
S^{2} &= x_{1}^{2}+2x_{1}x_{2}+x_{2}^{2} &&\justif{identité remarquable}\\
x_{1}^{2}+x_{2}^{2} &= S^{2}-2P &&\justif{on isole $x_{1}^{2}+x_{2}^{2}$ ; $x_{1}x_{2}=P$}\\
x_{1}^{2}+x_{2}^{2} &= \dfrac{81}{4}-\dfrac{4\fois{4}}{1\fois{4}} &&\justif{$S^{2}=\frac{81}{4}$ et $2P=4$}\\
x_{1}^{2}+x_{2}^{2} &= \res{\dfrac{65}{4}} &&
\end{calculs}
\rem{Contrôle (non demandé) : les racines sont $\frac{1}{2}$ et $4$ ;
$2+\frac{1}{4}=\frac{9}{4}$ et $\frac{1}{4}+16=\frac{65}{4}$.}
\end{exo}

\begin{exo}[Signe des racines]
\explique{On vérifie d'abord que $\Delta>0$. Si $P<0$, le produit de deux nombres est
négatif : ils sont de signes contraires. Si $P>0$, ils ont le même signe, et c'est
celui de leur somme $S$.}{règle des signes d'un produit}

\cas{1)}
\begin{calculs}
\Delta &= 81-56=25 &&\justif{$\Delta>0$ : deux racines}\\
P &= 14 &&\justif{$P>0$ : même signe}\\
S &= 9 &&\justif{$S>0$ : ce signe commun est $+$}
\end{calculs}
\explique{\res{Deux racines positives}.}{contrôle : ce sont $2$ et $7$}
\cas{2)}
\begin{calculs}
\Delta &= 4+60=64 &&\justif{$\Delta>0$}\\
P &= -15 &&\justif{$P<0$ : signes contraires}
\end{calculs}
\explique{\res{Deux racines de signes contraires}.}{contrôle : $-5$ et $3$ ; $S$ est
inutile ici}
\cas{3)}
\begin{calculs}
\Delta &= 100-36=64 &&\justif{$\Delta>0$}\\
P &= \dfrac{3}{3}=1 &&\justif{$P>0$ : même signe}\\
S &= -\dfrac{10}{3} &&\justif{$S<0$ : ce signe commun est $-$}
\end{calculs}
\explique{\res{Deux racines négatives}.}{contrôle : $-3$ et $-\frac{1}{3}$}
\end{exo}

\partie{Signe du trinôme et inéquations}

\begin{exo}[Tableau de signes]
\explique{On cherche les racines, puis on applique le théorème 5 : avec deux racines, le
trinôme a le signe de $a$ à l'extérieur, le signe contraire entre elles.}{sans racine :
toujours le signe de $a$}

\cas{A}
\begin{calculs}
2^{2}-6\times 2+8 &= 0 &&\justif{$x_{1}=2$ est racine évidente ($1$ ne convient pas)}\\
x_{2} &= 6-2=4 &&\justif{somme des racines $-\frac{b}{a}=6$}
\end{calculs}
\explique{$a=1>0$ : positif à l'extérieur de $2$ et $4$, négatif entre.}{}
\begin{center}
\begin{tikzpicture}[scale=0.8]
  \tkzTabInit[lgt=3,espcl=1.8]{$x$/0.8,$x^{2}-6x+8$/0.8}%
    {$-\infty$,$2$,$4$,$+\infty$}
  \tkzTabLine{,+,z,-,z,+,}
\end{tikzpicture}
\end{center}
\cas{B}
\begin{calculs}
-3+9-6 &= 0 &&\justif{$x_{1}=1$ : $a+b+c=0$}\\
S &= -\dfrac{9}{-3}=3 &&\\
x_{2} &= 3-1=2 &&
\end{calculs}
\explique{$a=-3<0$ : négatif à l'extérieur de $1$ et $2$, positif entre.}{}
\begin{center}
\begin{tikzpicture}[scale=0.8]
  \tkzTabInit[lgt=3,espcl=1.8]{$x$/0.8,$-3x^{2}+9x-6$/0.8}%
    {$-\infty$,$1$,$2$,$+\infty$}
  \tkzTabLine{,-,z,+,z,-,}
\end{tikzpicture}
\end{center}
\cas{C}
\begin{calculs}
\Delta &= (-4)^{2}-4\times 1\times 6 &&\\
\Delta &= 16-24=-8 &&\justif{$\Delta<0$ : aucune racine}
\end{calculs}
\explique{Le trinôme a toujours le signe de $a=1>0$.}{une seule case dans le tableau}
\begin{center}
\begin{tikzpicture}[scale=0.8]
  \tkzTabInit[lgt=3,espcl=3.6]{$x$/0.8,$x^{2}-4x+6$/0.8}%
    {$-\infty$,$+\infty$}
  \tkzTabLine{,+,}
\end{tikzpicture}
\end{center}
\cas{D}
\begin{calculs}
\Delta &= (-4)^{2}-4\times(-1)\times(-4) &&\\
\Delta &= 16-16=0 &&\justif{une racine double}\\
x_{0} &= -\dfrac{-4}{2\times(-1)}=-2 &&\justif{$x_{0}=-\frac{b}{2a}$}
\end{calculs}
\explique{$-x^{2}-4x-4=-(x+2)^{2}$ : du signe de $a=-1<0$, nul en $-2$.}{un carré
précédé d'un « $-$ » est négatif}
\begin{center}
\begin{tikzpicture}[scale=0.8]
  \tkzTabInit[lgt=3,espcl=1.8]{$x$/0.8,$-x^{2}-4x-4$/0.8}%
    {$-\infty$,$-2$,$+\infty$}
  \tkzTabLine{,-,z,-,}
\end{tikzpicture}
\end{center}
\end{exo}

\begin{exo}[Inéquations du second degré]
\explique{On cherche les racines, on en déduit le signe du trinôme (théorème 5), puis on
lit les $x$ qui conviennent. Crochets ouverts pour $<$ ou $>$, fermés aux racines pour
$\le$ ou $\ge$.}{méthode 6}

\cas{1)}
\explique{$x^{2}-7x+10>0$ : $a=1$, $b=-7$ et $c=10$.}{}
\begin{calculs}
\Delta &= 49-40=9=3^{2} &&\\
x_{1} &= \dfrac{7-3}{2}=2 &&\\
x_{2} &= \dfrac{7+3}{2}=5 &&
\end{calculs}
\begin{center}
\begin{tikzpicture}[scale=0.8]
  \tkzTabInit[lgt=3,espcl=1.8]{$x$/0.8,$x^{2}-7x+10$/0.8}%
    {$-\infty$,$2$,$5$,$+\infty$}
  \tkzTabLine{,+,z,-,z,+,}
\end{tikzpicture}
\end{center}
\explique{$S=\res{\left]-\infty\,;2\right[\cup\left]5\,;+\infty\right[}$.}{$a>0$ :
positif à l'extérieur ; $>$ strict, racines exclues}
\cas{2)}
\explique{$4$ est racine évidente ($16+8-24=0$) ; la somme des racines vaut $-2$, donc
l'autre est $-2-4=-6$.}{}
\explique{$S=\res{\left]-6\,;4\right[}$.}{$a=1>0$ : négatif entre les racines}
\cas{3)}
\explique{$p^{2}-p-30\le 0$ : $a=1$, $b=-1$ et $c=-30$.}{l'inconnue est $p$}
\begin{calculs}
\Delta &= 1+120=121=11^{2} &&\\
p_{1} &= \dfrac{1-11}{2}=-5 &&\\
p_{2} &= \dfrac{1+11}{2}=6 &&
\end{calculs}
\explique{$S=\res{\left[-5\,;6\right]}$.}{$a>0$ : négatif entre les racines ; $\le$ :
racines comprises}
\cas{4)}
\begin{calculs}
\Delta &= 3^{2}-4\times 1\times 4 &&\\
\Delta &= 9-16=-7 &&\justif{$\Delta<0$ : aucune racine}
\end{calculs}
\explique{Le trinôme a toujours le signe de $a=1>0$ : il n'est jamais négatif.
$S=\res{\varnothing}$.}{}
\cas{5)}
\begin{calculs}
2x^{2}-12x+18 &= 2\left(x^{2}-6x+9\right) &&\justif{facteur commun $2$}\\
2x^{2}-12x+18 &= 2(x-3)^{2} &&\justif{identité remarquable}
\end{calculs}
\explique{$2(x-3)^{2}$ est strictement positif, sauf en $3$ où il est nul.
$S=\res{\R\setminus\{3\}}$.}{$>0$ strict : on retire seulement $3$}
\cas{6)}
\explique{$-3x^{2}+2x+2>0$ : $a=-3$, $b=2$ et $c=2$.}{}
\begin{calculs}
\Delta &= 4+24=28 &&\justif{$\sqrt{28}=\sqrt{4\times 7}=2\sqrt{7}$}\\
x_{1} &= \dfrac{-2+2\sqrt{7}}{-6}=\dfrac{1-\sqrt{7}}{3} &&\justif{on divise haut et bas par $-2$}\\
x_{2} &= \dfrac{-2-2\sqrt{7}}{-6}=\dfrac{1+\sqrt{7}}{3} &&
\end{calculs}
\explique{$S=\res{\left]\dfrac{1-\sqrt{7}}{3}\,;\dfrac{1+\sqrt{7}}{3}\right[}$.}{$a<0$ :
positif entre les racines ; $x_{1}\approx -0{,}55$ et $x_{2}\approx 1{,}22$}
\rem{En 4) et 5), le trinôme ne change pas de signe : c'est le signe de $a$ qui décide.}
\end{exo}

\begin{exo}[Position d'une parabole et d'une droite]
\cas{1)}
\begin{calculs}
f(x)=4x+2 &\iff x^{2}+2x+3=4x+2 &&\justif{on remplace $f(x)$}\\
&\iff x^{2}-2x+1=0 &&\justif{tout dans le membre de gauche}\\
&\iff (x-1)^{2}=0 &&\justif{identité remarquable}\\
&\iff x=1 &&
\end{calculs}
\explique{L'équation a une seule solution : $\mathcal{C}_{f}$ et $(D)$ ont
\res{un seul point commun}, d'abscisse $1$ et d'ordonnée $4\times 1+2=6$, le point
$(1\,;6)$.}{contrôle : $f(1)=1+2+3=6$}
\cas{2)}
\begin{calculs}
f(x)-(4x+2) &= x^{2}+2x+3-4x-2 &&\justif{parenthèses : on retranche les deux termes}\\
f(x)-(4x+2) &= x^{2}-2x+1 &&\justif{on réduit}\\
f(x)-(4x+2) &= (x-1)^{2} &&\justif{identité remarquable}
\end{calculs}
\explique{Un carré est positif : $f(x)-(4x+2)$ est \res{positif ou nul}, nul seulement
en $1$.}{}
\cas{3)}
\explique{Pour tout réel $x$, $f(x)\ge 4x+2$ : \res{$\mathcal{C}_{f}$ est au-dessus de
$(D)$}, et elles se touchent au seul point $(1\,;6)$.}{signe de la
différence positif : la courbe de $f$ est au-dessus}
\rem{La droite touche la parabole en un point sans la traverser : c'est le cas
$\Delta=0$ de l'équation $f(x)=4x+2$.}
\end{exo}

\partie{Représentation graphique}

\begin{exo}[Sommet, racines, allure]
\explique{Pour chaque trinôme : signe de $a$ (orientation), $\alpha$ et $\beta$
(sommet), racines (points sur l'axe des abscisses), puis le point $(0\,;c)$ et son
symétrique par rapport à la droite $x=\alpha$.}{méthode 7}

\cas{A}
\begin{calculs}
a &= 1>0 &&\justif{parabole tournée vers le haut}\\
\alpha &= -\dfrac{4}{2}=-2 &&\justif{$\alpha=-\frac{b}{2a}$}\\
\beta &= 4-8+3 &&\justif{$\beta=A(-2)=(-2)^{2}+4\times(-2)+3$}\\
\beta &= -1 &&\justif{sommet $S(-2\,;-1)$}\\
\Delta &= 16-12=4 &&\justif{$\Delta>0$ : deux racines}\\
x_{1} &= \dfrac{-4-2}{2}=-3 &&\\
x_{2} &= \dfrac{-4+2}{2}=-1 &&
\end{calculs}
\explique{\res{Vers le haut}, sommet \res{$S(-2\,;-1)$}, racines \res{$-3$ et $-1$} ;
point $(0\,;3)$ et son symétrique $(-4\,;3)$.}{$\alpha=-2$ est bien le milieu de $-3$ et
$-1$}
\cas{B}
\begin{calculs}
a &= -1<0 &&\justif{tournée vers le bas}\\
\alpha &= -\dfrac{4}{-2}=2 &&\\
\beta &= -4+8-3 &&\justif{$\beta=B(2)$ ; $-2^{2}=-4$}\\
\beta &= 1 &&\justif{sommet $S(2\,;1)$}\\
\Delta &= 16-12=4 &&\justif{$\Delta=4^{2}-4\times(-1)\times(-3)$}\\
x_{1} &= \dfrac{-4+2}{-2}=1 &&\\
x_{2} &= \dfrac{-4-2}{-2}=3 &&
\end{calculs}
\explique{\res{Vers le bas}, sommet \res{$S(2\,;1)$}, racines \res{$1$ et $3$}.}{}
\cas{C}
\begin{calculs}
a &= 1>0 &&\justif{tournée vers le haut}\\
\alpha &= -\dfrac{-4}{2}=2 &&\\
\beta &= 4-8+5 &&\justif{$\beta=C(2)$}\\
\beta &= 1 &&\justif{sommet $S(2\,;1)$}\\
\Delta &= 16-20=-4 &&\justif{$\Delta<0$ : la courbe ne coupe pas l'axe}
\end{calculs}
\explique{\res{Vers le haut}, sommet \res{$S(2\,;1)$}, \res{aucune racine}. Pour
l'allure : $(0\,;5)$ et son symétrique $(4\,;5)$.}{sans racine, ces deux points
remplacent ceux de l'axe}
\cas{D}
\begin{calculs}
a &= -2<0 &&\justif{tournée vers le bas}\\
\alpha &= -\dfrac{-4}{2\times(-2)}=-1 &&\justif{$-\frac{-4}{-4}$ : résultat négatif}\\
\beta &= -2+4 &&\justif{$\beta=D(-1)=-2\times(-1)^{2}-4\times(-1)$}\\
\beta &= 2 &&\justif{sommet $S(-1\,;2)$}\\
-2x^{2}-4x &= -2x(x+2) &&\justif{pas de terme constant : facteur commun, pas de $\Delta$}
\end{calculs}
\explique{\res{Vers le bas}, sommet \res{$S(-1\,;2)$}, racines \res{$-2$ et $0$}.}{produit
nul}
\begin{center}
\begin{tikzpicture}
\begin{axis}[repere,x=0.42cm,y=0.42cm,
    xmin=-4.8,xmax=0.8,ymin=-1.8,ymax=3.8,
    xtick={-4,...,-1},ytick={-1,...,3},title={\scriptsize A},title style={yshift=-4pt}]
  \addplot[courbe,domain=-4.1:0.1]{x^2+4*x+3};
  \addplot[croix] coordinates {(-3,0)(-1,0)(-2,-1)(0,3)(-4,3)};
\end{axis}
\end{tikzpicture}\hfill
\begin{tikzpicture}
\begin{axis}[repere,x=0.42cm,y=0.42cm,
    xmin=-0.8,xmax=4.8,ymin=-3.8,ymax=1.8,
    xtick={1,...,4},ytick={-3,...,1},title={\scriptsize B},title style={yshift=-4pt}]
  \addplot[courbe,domain=-0.1:4.1]{-x^2+4*x-3};
  \addplot[croix] coordinates {(1,0)(3,0)(2,1)};
\end{axis}
\end{tikzpicture}\hfill
\begin{tikzpicture}
\begin{axis}[repere,x=0.42cm,y=0.42cm,
    xmin=-0.8,xmax=4.8,ymin=-0.8,ymax=6.3,
    xtick={1,...,4},ytick={1,...,6},title={\scriptsize C},title style={yshift=-4pt}]
  \addplot[courbe,domain=-0.2:4.2]{x^2-4*x+5};
  \addplot[croix] coordinates {(2,1)(0,5)(4,5)};
\end{axis}
\end{tikzpicture}\hfill
\begin{tikzpicture}
\begin{axis}[repere,x=0.42cm,y=0.42cm,
    xmin=-2.8,xmax=1.3,ymin=-3,ymax=2.8,
    xtick={-2,...,1},ytick={-2,...,2},title={\scriptsize D},title style={yshift=-4pt}]
  \addplot[courbe,domain=-2.55:0.55]{-2*x^2-4*x};
  \addplot[croix] coordinates {(-2,0)(0,0)(-1,2)};
\end{axis}
\end{tikzpicture}
\end{center}
\rem{On place d'abord le sommet et les racines, puis on trace une courbe « en U »
symétrique par rapport à la droite $x=\alpha$, arrondie au sommet (jamais pointue).}
\end{exo}

\begin{exo}[Utiliser la symétrie]
\cas{1)}
\explique{$f(x)=x^{2}-6x+8$ : $a=1$, $b=-6$ et $c=8$.}{}
\begin{calculs}
\alpha &= -\dfrac{-6}{2}=\res{3} &&\justif{$\alpha=-\frac{b}{2a}$}\\
\beta &= f(3)=9-18+8 &&\\
\beta &= \res{-1} &&\\
2^{2}-6\times 2+8 &= 0 &&\justif{$2$ est racine évidente}\\
x_{2} &= 6-2=4 &&\justif{somme des racines $6$}
\end{calculs}
\explique{Les racines sont \res{$2$ et $4$}.}{contrôle : $\alpha=3$ est leur milieu}
\cas{2)}
\begin{calculs}
f(0) &= \res{8} &&\justif{forme développée : le terme constant}
\end{calculs}
\explique{La parabole est symétrique par rapport à la droite $x=3$. Le point $(0\,;8)$
est à $3$ unités à gauche de cette droite ; son symétrique est à $3$ unités à droite,
à la même hauteur : abscisse $3+3=6$, c'est \res{$(6\,;8)$}.}{la symétrie garde
l'ordonnée}
\begin{calculs}
f(6) &= 36-36+8 &&\justif{vérification par le calcul}\\
f(6) &= 8 &&\justif{$(6\,;8)$ est bien sur la parabole}
\end{calculs}
\cas{3)}
\begin{minipage}[c]{0.56\linewidth}\raggedright
\explique{On place le sommet $S(3\,;-1)$, les points $(2\,;0)$ et $(4\,;0)$, puis
$(0\,;8)$ et $(6\,;8)$, et on trace la parabole tournée vers le haut ($a>0$).}{la droite
$x=3$, en pointillés, aide à respecter la symétrie}
\end{minipage}\hfill
\begin{minipage}[c]{0.4\linewidth}\centering
\begin{tikzpicture}
\begin{axis}[repere,x=0.38cm,y=0.38cm,
    xmin=-0.8,xmax=7.5,ymin=-1.8,ymax=9.3,
    xtick={1,...,7},ytick={-1,...,8}]
  \addplot[courbe,domain=-0.1:6.1]{x^2-6*x+8};
  \addplot[croix] coordinates {(2,0)(4,0)(3,-1)(0,8)(6,8)};
  \addplot[dashed,black!50] coordinates {(3,-1.8)(3,9.3)};
\end{axis}
\end{tikzpicture}
\end{minipage}
\end{exo}

\begin{exo}[Retrouver un trinôme]
\cas{1)}
\begin{calculs}
\dfrac{-1+3}{2} &= 1 &&\justif{milieu des racines $-1$ et $3$}
\end{calculs}
\explique{C'est l'abscisse du sommet $S(1\,;-8)$ : \res{l'abscisse du sommet est bien le
milieu des racines}.}{symétrie de la parabole}
\cas{2)}
\explique{Le sommet est un point de la parabole : en $x=1$, le trinôme vaut $-8$.}{on
traduit « $S$ est sur la courbe » par une égalité}
\begin{calculs}
a(1+1)(1-3)=-8 &\iff a\times 2\times(-2)=-8 &&\justif{on remplace $x$ par $1$}\\
&\iff -4a=-8 &&\\
&\iff a=\res{2} &&\justif{on divise par $-4$}
\end{calculs}
\cas{3)}
\begin{calculs}
2(x+1)(x-3) &= 2\left(x^{2}-3x+x-3\right) &&\justif{double distributivité, entre parenthèses}\\
2(x+1)(x-3) &= 2\left(x^{2}-2x-3\right) &&\justif{on réduit}\\
2(x+1)(x-3) &= \res{2x^{2}-4x-6} &&\justif{on distribue le $2$}
\end{calculs}
\explique{Contrôle avec le sommet : $\alpha=-\dfrac{-4}{4}=1$ et $2-4-6=-8$, on
retrouve $S(1\,;-8)$.}{}
\end{exo}

\begin{exo}[Le jet d'une fontaine]
\cas{1)}
\explique{Les racines de $P$ sont les abscisses des points où la parabole coupe l'axe
des abscisses : $x_{1}=\res{1}$ et $x_{2}=\res{5}$.}{on lit sur l'axe horizontal}
\explique{À ces distances, la hauteur du jet est nulle : le jet sort du bassin à $1$ m
de l'origine et y retombe à $5$ m.}{on interprète dans le contexte}
\cas{2)}
\explique{Le sommet est le point le plus haut de la parabole :
$(\alpha\,;\beta)=(\res{3}\,;\res{2})$.}{lecture sur le quadrillage}
\explique{$\beta$ est la \res{hauteur maximale du jet} : $2$ m, atteinte à $3$ m de
l'origine.}{$\beta$ est le maximum du trinôme}
\cas{3)}
\begin{calculs}
\dfrac{x_{1}+x_{2}}{2} &= \dfrac{1+5}{2} &&\\
\dfrac{x_{1}+x_{2}}{2} &= 3 &&\justif{c'est bien $\alpha$}
\end{calculs}
\cas{4)}
\explique{La parabole est tournée vers le bas (le trinôme a un maximum) : $a$ est
\res{strictement négatif}.}{théorème 6 : $a>0$ donnerait une parabole tournée vers le
haut}
\end{exo}

\begin{exo}[À chaque parabole sa fonction]
\explique{On utilise le signe de $a$ (orientation), le signe de $\Delta$ (nombre de
points sur l'axe des abscisses), puis les racines ou le sommet pour départager.}{trois
paraboles passent par $(0\,;2)$ : $c$ seul ne suffit pas}

\cas{a)}
\begin{calculs}
a &= -1<0 &&\justif{tournée vers le bas}\\
\Delta &= 4-8=-4 &&\justif{$(-2)^{2}-4\times(-1)\times(-2)$ : aucune racine}\\
\alpha &= -\dfrac{-2}{-2}=-1 &&\\
\beta &= -1+2-2=-1 &&\justif{sommet $(-1\,;-1)$}
\end{calculs}
\explique{Vers le bas et entièrement sous l'axe : $f_{1}$ a pour courbe
\res{$\mathcal{C}_{3}$}.}{$\mathcal{C}_{1}$ est aussi vers le bas, mais coupe l'axe}
\cas{b)}
\begin{calculs}
a &= 1>0 &&\justif{tournée vers le haut}\\
\Delta &= 4-8=-4 &&\justif{aucune racine}\\
\alpha &= 1 &&\\
\beta &= 1-2+2=1 &&\justif{sommet $(1\,;1)$}
\end{calculs}
\explique{Vers le haut et entièrement au-dessus de l'axe : \res{$\mathcal{C}_{2}$}.}{}
\cas{c)}
\begin{calculs}
a &= 2>0 &&\justif{tournée vers le haut}\\
\Delta &= 25-16=9 &&\justif{deux racines}\\
x_{1} &= \dfrac{-5-3}{4}=-2 &&\\
x_{2} &= \dfrac{-5+3}{4}=-\dfrac{1}{2} &&
\end{calculs}
\explique{Vers le haut, coupe l'axe en $-2$ et en $-\dfrac{1}{2}$ :
\res{$\mathcal{C}_{5}$}.}{}
\cas{d)}
\begin{calculs}
a &= -2<0 &&\justif{tournée vers le bas}\\
\Delta &= 9+16=25 &&\justif{deux racines}\\
x_{1} &= \dfrac{-3+5}{-4}=-\dfrac{1}{2} &&\\
x_{2} &= \dfrac{-3-5}{-4}=2 &&
\end{calculs}
\explique{Vers le bas, coupe l'axe en $-\dfrac{1}{2}$ et en $2$ :
\res{$\mathcal{C}_{1}$}.}{}
\cas{e)}
\begin{calculs}
\Delta &= 1-4\times\dfrac{1}{4}=0 &&\justif{une racine double}\\
x_{0} &= \dfrac{1}{2} &&\justif{$f_{5}(x)=\left(x-\frac{1}{2}\right)^{2}$}
\end{calculs}
\explique{Vers le haut ($a=1$), la courbe touche l'axe en $\dfrac{1}{2}$ seulement :
\res{$\mathcal{C}_{4}$}.}{le sommet est sur l'axe}
\rem{Contrôle : chaque parabole est attribuée une seule fois. On peut aussi vérifier
que $f_{2}(0)$, $f_{3}(0)$ et $f_{4}(0)$ valent $2$ sur la figure.}
\end{exo}

\partie{Variations de la fonction trinôme}

\begin{exo}[Variations et forme canonique]
\explique{Sous forme canonique, on lit $a$, $\alpha$ et $\beta$. Si $a>0$ : décroissante
puis croissante, minimum $\beta$ en $\alpha$ ; si $a<0$ : croissante puis
décroissante, maximum $\beta$ en $\alpha$.}{théorème 7}

\cas{1)}
\explique{$f(x)=3(x-2)^{2}+1$ : $a=3>0$, $\alpha=2$ et $\beta=1$.}{lecture directe}
\begin{center}
\begin{tikzpicture}[scale=0.75]
  \tkzTabInit[lgt=1.4,espcl=2]{$x$/0.8,$f(x)$/1.4}{$-\infty$,$2$,$+\infty$}
  \tkzTabVar{+/{},-/$1$,+/{}}
\end{tikzpicture}
\end{center}
\explique{$f$ a un minimum, égal à \res{$1$}, atteint en $2$.}{un « creux » car $a>0$}
\cas{2)}
\explique{$g(x)=-2(x+3)^{2}+4=-2\bigl(x-(-3)\bigr)^{2}+4$ : $a=-2<0$, $\alpha=-3$ et
$\beta=4$.}{$x+3=x-(-3)$}
\begin{center}
\begin{tikzpicture}[scale=0.75]
  \tkzTabInit[lgt=1.4,espcl=2]{$x$/0.8,$g(x)$/1.4}{$-\infty$,$-3$,$+\infty$}
  \tkzTabVar{-/{},+/$4$,-/{}}
\end{tikzpicture}
\end{center}
\explique{$g$ a un maximum, égal à \res{$4$}, atteint en $-3$.}{une « bosse » car $a<0$}
\cas{3)}
\explique{$h(x)=x(x+6)$ est une forme factorisée : ses racines $0$ et $-6$ se
lisent.}{produit nul}
\begin{calculs}
\alpha &= \dfrac{-6+0}{2}=-3 &&\justif{$\alpha$ est le milieu des racines}\\
\beta &= h(-3)=-3\times 3 &&\justif{$-3\times(-3+6)$}\\
\beta &= -9 &&\\
h(x) &= x^{2}+6x &&\justif{en développant : $a=1>0$}
\end{calculs}
\begin{center}
\begin{tikzpicture}[scale=0.75]
  \tkzTabInit[lgt=1.4,espcl=2]{$x$/0.8,$h(x)$/1.4}{$-\infty$,$-3$,$+\infty$}
  \tkzTabVar{+/{},-/$-9$,+/{}}
\end{tikzpicture}
\end{center}
\explique{$h$ a un minimum, égal à \res{$-9$}, atteint en $-3$.}{}
\rem{Dans $a(x-\alpha)^{2}+\beta$, attention au signe : $(x+3)^{2}=\bigl(x-(-3)\bigr)^{2}$
donne $\alpha=-3$.}
\end{exo}

\begin{exo}[Variations par $\alpha$ et $\beta$]
\cas{A}
\begin{calculs}
a &= 1>0 &&\justif{décroissante, puis croissante}\\
\alpha &= -\dfrac{6}{2}=-3 &&\justif{$\alpha=-\frac{b}{2a}$}\\
\beta &= 9-18+7 &&\justif{$\beta=A(-3)$}\\
\beta &= -2 &&
\end{calculs}
\begin{center}
\begin{tikzpicture}[scale=0.75]
  \tkzTabInit[lgt=1.2,espcl=1.8]{$x$/0.8,$A$/1.4}{$-\infty$,$-3$,$+\infty$}
  \tkzTabVar{+/{},-/$-2$,+/{}}
\end{tikzpicture}
\end{center}
\explique{\res{Décroissante puis croissante}, minimum $-2$ en $-3$.}{}
\cas{B}
\begin{calculs}
a &= -1<0 &&\justif{croissante, puis décroissante}\\
\alpha &= -\dfrac{4}{-2}=2 &&\\
\beta &= -4+8+5 &&\justif{$\beta=B(2)$ ; $-2^{2}=-4$}\\
\beta &= 9 &&
\end{calculs}
\begin{center}
\begin{tikzpicture}[scale=0.75]
  \tkzTabInit[lgt=1.2,espcl=1.8]{$x$/0.8,$B$/1.4}{$-\infty$,$2$,$+\infty$}
  \tkzTabVar{-/{},+/$9$,-/{}}
\end{tikzpicture}
\end{center}
\explique{\res{Croissante puis décroissante}, maximum $9$ en $2$.}{}
\cas{C}
\begin{calculs}
a &= 3>0 &&\\
\alpha &= -\dfrac{-12}{6}=2 &&\\
\beta &= 12-24+5 &&\justif{$\beta=C(2)=3\times 4-12\times 2+5$}\\
\beta &= -7 &&
\end{calculs}
\begin{center}
\begin{tikzpicture}[scale=0.75]
  \tkzTabInit[lgt=1.2,espcl=1.8]{$x$/0.8,$C$/1.4}{$-\infty$,$2$,$+\infty$}
  \tkzTabVar{+/{},-/$-7$,+/{}}
\end{tikzpicture}
\end{center}
\explique{\res{Décroissante puis croissante}, minimum $-7$ en $2$.}{}
\cas{D}
\explique{$D=-\dfrac{1}{2}x^{2}-2x$ : $a=-\dfrac{1}{2}$, $b=-2$ et $c=0$.}{pas de terme
constant : $c=0$}
\begin{calculs}
\alpha &= -\dfrac{-2}{2\times\left(-\frac{1}{2}\right)} &&\justif{$2a=-1$}\\
\alpha &= -\dfrac{-2}{-1}=-2 &&\\
\beta &= -\dfrac{1}{2}\times 4+4 &&\justif{$\beta=D(-2)$ ; $-2\times(-2)=4$}\\
\beta &= 2 &&
\end{calculs}
\begin{center}
\begin{tikzpicture}[scale=0.75]
  \tkzTabInit[lgt=1.2,espcl=1.8]{$x$/0.8,$D$/1.4}{$-\infty$,$-2$,$+\infty$}
  \tkzTabVar{-/{},+/$2$,-/{}}
\end{tikzpicture}
\end{center}
\explique{\res{Croissante puis décroissante}, maximum $2$ en $-2$.}{$a<0$}
\end{exo}

\begin{exo}[Un maximum]
\cas{1)}
\begin{calculs}
f(x) &= -3x^{2}+12x-10 &&\justif{on recopie}\\
f(x) &= -3\left(x^{2}-4x\right)-10 &&\justif{$-3$ en facteur : $12x=-3\times(-4x)$}\\
f(x) &= -3\left[(x-2)^{2}-4\right]-10 &&\justif{$x^{2}-4x$ est le début de $(x-2)^{2}$}\\
f(x) &= -3(x-2)^{2}+12-10 &&\justif{on distribue le $-3$}\\
f(x) &= \res{-3(x-2)^{2}+2} &&\justif{on réduit}
\end{calculs}
\explique{Contrôle : $\alpha=-\dfrac{12}{-6}=2$ et $f(2)=-12+24-10$, soit $2$ : c'est $\beta$.}{méthode 2}
\cas{2)}
\explique{On part de la forme canonique et d'un fait sûr : un carré est
positif.}{démonstration en colonne}
\begin{calculs}
(x-2)^{2} &\ge 0 &&\justif{un carré est positif}\\
-3(x-2)^{2} &\le 0 &&\justif{on multiplie par $-3<0$ : le sens change}\\
-3(x-2)^{2}+2 &\le 2 &&\justif{on ajoute $2$}\\
f(x) &\le 2 &&\justif{pour tout réel $x$}
\end{calculs}
\begin{calculs}
f(x)=2 &\iff -3(x-2)^{2}=0 &&\\
&\iff x=2 &&\justif{un carré est nul seulement si le nombre est nul}
\end{calculs}
\explique{Le maximum de $f$ est \res{$2$}, atteint pour \res{$x=2$}.}{}
\rem{Erreur fréquente : garder le sens de l'inégalité en multipliant par $-3$.}
\end{exo}

\begin{exo}[Étude complète d'un trinôme]
\cas{1)}
\explique{$f(x)=-x^{2}+2x+3$ : $a=-1$, $b=2$ et $c=3$.}{}
\begin{calculs}
-(-1)^{2}+2\times(-1)+3 &= 0 &&\justif{$x_{1}=-1$ est racine évidente ($1$ donne $4$)}\\
S &= -\dfrac{2}{-1}=2 &&\justif{$S=-\frac{b}{a}$}\\
x_{2} &= 2-(-1)=3 &&
\end{calculs}
\explique{Les racines sont \res{$-1$ et $3$}, donc
$f(x)=-1\times(x+1)(x-3)=\res{-(x+1)(x-3)}$.}{$a=-1$ devant ; contrôle :
$-1\times(-1)\times 3=3=c$}
\cas{2)}
\begin{calculs}
\alpha &= -\dfrac{2}{2\times(-1)}=\res{1} &&\\
\beta &= f(1)=-1+2+3 &&\\
\beta &= \res{4} &&
\end{calculs}
\explique{$a=-1<0$ : croissante puis décroissante, maximum $4$ en $1$.}{théorème 7}
\begin{center}
\begin{tikzpicture}[scale=0.75]
  \tkzTabInit[lgt=1.4,espcl=2]{$x$/0.8,$f(x)$/1.4}{$-\infty$,$1$,$+\infty$}
  \tkzTabVar{-/{},+/$4$,-/{}}
\end{tikzpicture}
\end{center}
\cas{3)}
\begin{calculs}
\dfrac{-1+3}{2} &= 1 &&\justif{c'est bien $\alpha$}
\end{calculs}
\cas{4)}
\explique{$a<0$ : $f(x)$ est négatif à l'extérieur des racines, positif entre elles.
Pour la courbe : sommet $(1\,;4)$, racines, puis $(0\,;3)$ et son symétrique
$(2\,;3)$.}{théorème 5, puis méthode 7}
\begin{minipage}[t]{0.5\linewidth}
  \vspace{-4pt}
  \begin{tikzpicture}[scale=0.8]
    \tkzTabInit[lgt=1.6,espcl=1.6]{$x$/0.8,$f(x)$/0.8}%
      {$-\infty$,$-1$,$3$,$+\infty$}
    \tkzTabLine{,-,z,+,z,-,}
  \end{tikzpicture}
\end{minipage}\hfill
\begin{minipage}[t]{0.44\linewidth}\centering
  \vspace{-4pt}
  \begin{tikzpicture}
  \begin{axis}[repere,x=0.45cm,y=0.45cm,
      xmin=-1.8,xmax=3.8,ymin=-2.2,ymax=4.8,
      xtick={-1,...,3},ytick={-2,...,4}]
    \addplot[courbe,domain=-1.4:3.4]{-x^2+2*x+3};
    \addplot[croix] coordinates {(-1,0)(3,0)(1,4)(0,3)(2,3)};
  \end{axis}
  \end{tikzpicture}
\end{minipage}
\rem{Le tableau de signes se lit sur la courbe : elle est au-dessus de l'axe entre $-1$
et $3$, en dessous ailleurs.}
\end{exo}

\partie{Équations se ramenant au second degré}

\begin{exo}[Développer, puis résoudre]
\explique{On développe, on passe tout dans un membre pour obtenir $ax^{2}+bx+c=0$, puis
on résout (factorisation directe, racine évidente ou discriminant).}{méthode 9}

\cas{1)}
\begin{calculs}
(x+1)(x-4)=6 &\iff x^{2}-4x+x-4=6 &&\justif{on développe}\\
&\iff x^{2}-3x-10=0 &&\justif{tout dans un membre}
\end{calculs}
\begin{calculs}
\Delta &= 9+40=49=7^{2} &&\justif{$a=1$, $b=-3$ et $c=-10$}\\
x_{1} &= \dfrac{3+7}{2}=5 &&\\
x_{2} &= \dfrac{3-7}{2}=-2 &&
\end{calculs}
\explique{$S=\res{\{-2\,;5\}}$.}{}
\rem{Erreur fréquente : écrire « $x+1=6$ ou $x-4=6$ ». La règle du produit nul ne vaut
que si le produit est égal à $0$.}
\cas{2)}
\begin{calculs}
(3x+1)^{2}=16 &\iff (3x+1)^{2}-4^{2}=0 &&\justif{tout dans un membre}\\
&\iff (3x+1-4)(3x+1+4)=0 &&\justif{$u^{2}-v^{2}$}\\
&\iff 3x-3=0 \text{ ou } 3x+5=0 &&\justif{produit nul}\\
&\iff x=1 \text{ ou } x=-\dfrac{5}{3} &&
\end{calculs}
\explique{$S=\res{\left\{-\dfrac{5}{3}\,;1\right\}}$.}{}
\rem{Plus rapide : $A^{2}=16 \iff A=4$ ou $A=-4$. On n'oublie pas la valeur
négative.}
\cas{3)}
\begin{calculs}
x(x-4)=-3 &\iff x^{2}-4x+3=0 &&\justif{on développe ; tout dans un membre}
\end{calculs}
\explique{$1$ est racine évidente ($1-4+3=0$) ; la somme des racines vaut $4$, donc
l'autre est $3$. $S=\res{\{1\,;3\}}$.}{méthode 5}
\cas{4)}
\begin{calculs}
(2x+3)^{2}=(x+3)^{2} &\iff 4x^{2}+12x+9=x^{2}+6x+9 &&\justif{on développe}\\
&\iff 3x^{2}+6x=0 &&\justif{tout dans un membre : les $9$ disparaissent}\\
&\iff 3x(x+2)=0 &&\justif{facteur commun : pas de $\Delta$}
\end{calculs}
\explique{$S=\res{\{-2\,;0\}}$.}{produit nul}
\cas{5)}
\begin{calculs}
(x-3)^{2}=x+3 &\iff x^{2}-6x+9=x+3 &&\justif{on développe}\\
&\iff x^{2}-7x+6=0 &&
\end{calculs}
\explique{$1$ est racine évidente ($1-7+6=0$) ; l'autre est $7-1=6$.
$S=\res{\{1\,;6\}}$.}{somme des racines $7$}
\cas{6)}
\begin{calculs}
(x+4)(x-4)=6x+11 &\iff x^{2}-16=6x+11 &&\justif{$(u+v)(u-v)=u^{2}-v^{2}$}\\
&\iff x^{2}-6x-27=0 &&
\end{calculs}
\begin{calculs}
\Delta &= 36+108=144=12^{2} &&\\
x_{1} &= \dfrac{6+12}{2}=9 &&\\
x_{2} &= \dfrac{6-12}{2}=-3 &&
\end{calculs}
\explique{$S=\res{\{-3\,;9\}}$.}{}
\end{exo}

\begin{exo}[Un seul quotient]
\explique{On écrit d'abord l'ensemble de définition $D$, on multiplie par le
dénominateur (non nul sur $D$), on résout, puis on ne garde que les solutions qui sont
dans $D$.}{méthode 9, exemple c)}

\cas{1)}
\explique{$x\ne 0$ : $D=\R\setminus\{0\}$.}{}
\begin{calculs}
\dfrac{12}{x}=x+1 &\iff 12=x(x+1) &&\justif{on multiplie par $x\ne 0$}\\
&\iff x^{2}+x-12=0 &&
\end{calculs}
\explique{$3$ est racine évidente ($9+3-12=0$) ; l'autre est $-1-3=-4$. Les deux sont
dans $D$ : $S=\res{\{-4\,;3\}}$.}{contrôle : $\frac{12}{3}=4=3+1$}
\cas{2)}
\explique{$x-1\ne 0$ : $D=\R\setminus\{1\}$.}{}
\begin{calculs}
\dfrac{6}{x-1}=x &\iff 6=x(x-1) &&\justif{on multiplie par $x-1\ne 0$}\\
&\iff x^{2}-x-6=0 &&
\end{calculs}
\explique{$3$ est racine évidente ($9-3-6=0$) ; l'autre est $1-3=-2$. Les deux sont
différentes de $1$ : $S=\res{\{-2\,;3\}}$.}{}
\cas{3)}
\explique{$D=\R\setminus\{0\}$.}{}
\begin{calculs}
\dfrac{x+10}{x}=x-2 &\iff x+10=x(x-2) &&\justif{on multiplie par $x\ne 0$}\\
&\iff x+10=x^{2}-2x &&\\
&\iff x^{2}-3x-10=0 &&
\end{calculs}
\explique{$5$ est racine évidente ($25-15-10=0$) ; l'autre est $3-5=-2$.
$S=\res{\{-2\,;5\}}$.}{toutes deux dans $D$}
\cas{4)}
\explique{$D=\R\setminus\{0\}$.}{}
\begin{calculs}
x+\dfrac{8}{x}=6 &\iff x^{2}+8=6x &&\justif{on multiplie \emph{chaque terme} par $x$}\\
&\iff x^{2}-6x+8=0 &&
\end{calculs}
\explique{$2$ est racine évidente ($4-12+8=0$) ; l'autre est $6-2=4$.
$S=\res{\{2\,;4\}}$.}{}
\rem{Erreur fréquente en 4) : oublier de multiplier le $x$ isolé, ce qui donne
$x+8=6x$.}
\end{exo}

\begin{exo}[Équations avec quotients]
\cas{1)}
\explique{$D=\R\setminus\{3\}$.}{}
\begin{calculs}
\dfrac{x^{2}-6x+9}{x-3}=2x-5 &\iff x^{2}-6x+9=(2x-5)(x-3) &&\justif{on multiplie par $x-3\ne 0$}\\
&\iff x^{2}-6x+9=2x^{2}-11x+15 &&\justif{on développe}\\
&\iff x^{2}-5x+6=0 &&\justif{tout dans le membre de droite}
\end{calculs}
\explique{$2$ est racine évidente ($4-10+6=0$) ; l'autre est $5-2=3$. \textbf{Mais}
$3\notin D$ : $S=\res{\{2\}}$.}{la valeur interdite est rejetée}
\rem{Plus rapide : $x^{2}-6x+9=(x-3)^{2}$, donc sur $D$ le quotient vaut $x-3$, et
$x-3=2x-5 \iff x=2$.}
\cas{2)}
\explique{$D=\R\setminus\{0\,;3\}$. On multiplie par $x(x-3)$, non nul sur $D$.}{le
dénominateur commun}
\begin{calculs}
\dfrac{2}{x}+\dfrac{2}{x-3}=1 &\iff 2(x-3)+2x=x(x-3) &&\justif{chaque terme multiplié par $x(x-3)$}\\
&\iff 4x-6=x^{2}-3x &&\\
&\iff x^{2}-7x+6=0 &&
\end{calculs}
\explique{Racines $1$ (évidente) et $7-1=6$, toutes deux dans $D$ :
$S=\res{\{1\,;6\}}$.}{contrôle en $1$ : $2+\frac{2}{-2}=1$}
\cas{3)}
\explique{$D=\R\setminus\{-2\,;3\}$. On multiplie par $2(x+2)(x-3)$, non nul sur
$D$.}{le $2$ vient du dénominateur de $\frac{3}{2}$}
\begin{calculs}
\dfrac{1}{x+2}-\dfrac{2}{x-3}=\dfrac{3}{2} &\iff 2(x-3)-4(x+2)=3(x+2)(x-3) &&\\
&\iff -2x-14=3x^{2}-3x-18 &&\justif{on développe les deux membres}\\
&\iff 3x^{2}-x-4=0 &&
\end{calculs}
\begin{calculs}
\Delta &= 1+48=49=7^{2} &&\\
x_{1} &= \dfrac{1+7}{6}=\dfrac{4}{3} &&\\
x_{2} &= \dfrac{1-7}{6}=-1 &&
\end{calculs}
\explique{Les deux sont dans $D$ : $S=\res{\left\{-1\,;\dfrac{4}{3}\right\}}$.}{contrôle
en $-1$ : $1-\frac{2}{-4}=\frac{3}{2}$}
\cas{4)}
\explique{$D=\R\setminus\{-2\}$.}{}
\begin{calculs}
\dfrac{2x^{2}+9x+10}{x+2}=x+3 &\iff 2x^{2}+9x+10=(x+3)(x+2) &&\justif{on multiplie par $x+2\ne 0$}\\
&\iff 2x^{2}+9x+10=x^{2}+5x+6 &&\\
&\iff x^{2}+4x+4=0 &&\\
&\iff (x+2)^{2}=0 &&\justif{identité remarquable}
\end{calculs}
\explique{La seule valeur est $-2$, qui n'est pas dans $D$ : $S=\res{\varnothing}$.}{on
vérifie toujours les solutions avec $D$}
\end{exo}

\partie{Inéquations se ramenant au second degré}

\begin{exo}[Développer, puis résoudre]
\explique{On développe, on passe tout dans un membre pour comparer à $0$, puis on étudie
le signe du trinôme.}{méthode 10, exemple a)}

\cas{1)}
\begin{calculs}
(x+1)(x+3)\le 2x+6 &\iff x^{2}+4x+3\le 2x+6 &&\justif{on développe}\\
&\iff x^{2}+2x-3\le 0 &&\justif{tout dans un membre}
\end{calculs}
\explique{$1$ est racine évidente ($1+2-3=0$) ; l'autre est $-2-1=-3$. $a=1>0$ : négatif
entre les racines. $S=\res{\left[-3\,;1\right]}$.}{$\le$ : racines comprises}
\rem{Autre voie : $2x+6=2(x+3)$, donc l'inéquation s'écrit $(x+3)(x+1-2)\le 0$,
soit $(x+3)(x-1)\le 0$.}
\cas{2)}
\begin{calculs}
(3x+1)^{2}<(x+3)^{2} &\iff 9x^{2}+6x+1<x^{2}+6x+9 &&\justif{on développe}\\
&\iff 8x^{2}-8<0 &&\\
&\iff 8(x-1)(x+1)<0 &&\justif{facteur $8$, puis $u^{2}-v^{2}$}
\end{calculs}
\explique{$S=\res{\left]-1\,;1\right[}$.}{$a=8>0$ : négatif entre $-1$ et $1$}
\cas{3)}
\begin{calculs}
x(x-4)\ge 5 &\iff x^{2}-4x-5\ge 0 &&\justif{on développe ; tout dans un membre}
\end{calculs}
\explique{$-1$ est racine évidente ($1+4-5=0$) ; l'autre est $4-(-1)=5$.
$S=\res{\left]-\infty\,;-1\right]\cup\left[5\,;+\infty\right[}$.}{$a>0$ : positif à
l'extérieur ; racines comprises}
\cas{4)}
\begin{calculs}
(x+2)^{2}>9 &\iff x^{2}+4x+4>9 &&\justif{on développe}\\
&\iff x^{2}+4x-5>0 &&
\end{calculs}
\explique{$1$ est racine évidente ; l'autre est $-4-1=-5$.
$S=\res{\left]-\infty\,;-5\right[\cup\left]1\,;+\infty\right[}$.}{$a>0$ : positif à
l'extérieur ; strict}
\rem{Erreur fréquente : passer de $(x+2)^{2}>9$ à $x+2>3$ seulement ; on oublie
$x+2<-3$, c'est-à-dire $x<-5$.}
\end{exo}

\begin{exo}[Signe d'un quotient]
\explique{On ne multiplie pas par le dénominateur (son signe change) : on étudie le signe
de chaque facteur dans un tableau, avec une double barre aux valeurs
interdites.}{méthode 10, exemple b)}

\cas{1)}
\explique{$x+1\ne 0$ : $D=\R\setminus\{-1\}$. $2x-6$ s'annule en $3$, $x+1$ en
$-1$.}{deux facteurs du premier degré}
\begin{center}
\begin{tikzpicture}[scale=0.8]
  \tkzTabInit[lgt=2,espcl=1.8]{$x$/0.8,$2x-6$/0.8,$x+1$/0.8,$\frac{2x-6}{x+1}$/1.1}%
    {$-\infty$,$-1$,$3$,$+\infty$}
  \tkzTabLine{,-,t,-,z,+,}
  \tkzTabLine{,-,z,+,t,+,}
  \tkzTabLine{,+,d,-,z,+,}
\end{tikzpicture}
\end{center}
\explique{$S=\res{\left]-1\,;3\right]}$.}{$3$ compris (quotient nul) ; $-1$ exclu
(valeur interdite)}
\cas{2)}
\explique{$D=\R\setminus\{-2\}$. $x^{2}-1=(x-1)(x+1)$ a pour racines $-1$ et $1$, et
$a=1>0$ : positif à l'extérieur.}{}
\begin{center}
\begin{tikzpicture}[scale=0.8]
  \tkzTabInit[lgt=2,espcl=1.8]{$x$/0.8,$x^{2}-1$/0.8,$x+2$/0.8,$\frac{x^{2}-1}{x+2}$/1.1}%
    {$-\infty$,$-2$,$-1$,$1$,$+\infty$}
  \tkzTabLine{,+,t,+,z,-,z,+,}
  \tkzTabLine{,-,z,+,t,+,t,+,}
  \tkzTabLine{,-,d,+,z,-,z,+,}
\end{tikzpicture}
\end{center}
\explique{$S=\res{\left]-2\,;-1\right]\cup\left[1\,;+\infty\right[}$.}{$\ge 0$ : on
prend les « $+$ » et les zéros du numérateur}
\cas{3)}
\explique{Dénominateur $x^{2}-4x+3$ : racines $1$ (évidente) et $4-1=3$, donc
$D=\R\setminus\{1\,;3\}$.}{somme des racines $4$}
\explique{Numérateur $x^{2}+x-6$ : racines $2$ (évidente : $4+2-6=0$) et $-1-2=-3$.}{somme
des racines $-1$}
\explique{Les deux trinômes ont $a=1>0$ : positifs à l'extérieur de leurs
racines.}{théorème 5}
\begin{center}
\begin{tikzpicture}[scale=0.8]
  \tkzTabInit[lgt=2.6,espcl=1.6]{$x$/0.8,$x^{2}+x-6$/0.8,$x^{2}-4x+3$/0.8,%
    $\frac{x^{2}+x-6}{x^{2}-4x+3}$/1.1}%
    {$-\infty$,$-3$,$1$,$2$,$3$,$+\infty$}
  \tkzTabLine{,+,z,-,t,-,z,+,t,+,}
  \tkzTabLine{,+,t,+,z,-,t,-,z,+,}
  \tkzTabLine{,+,z,-,d,+,z,-,d,+,}
\end{tikzpicture}
\end{center}
\explique{$S=\res{\left]-\infty\,;-3\right[\cup\left]1\,;2\right[\cup\left]3\,;+\infty\right[}$.}{$>0$
strict : zéros exclus}
\cas{4)}
\explique{$D=\R\setminus\{2\}$. Numérateur $x^{2}-5x+4$ : racines $1$ (évidente) et
$5-1=4$, positif à l'extérieur.}{}
\begin{center}
\begin{tikzpicture}[scale=0.8]
  \tkzTabInit[lgt=2,espcl=1.8]{$x$/0.8,$x^{2}-5x+4$/0.8,$x-2$/0.8,$\frac{x^{2}-5x+4}{x-2}$/1.1}%
    {$-\infty$,$1$,$2$,$4$,$+\infty$}
  \tkzTabLine{,+,z,-,t,-,z,+,}
  \tkzTabLine{,-,t,-,z,+,t,+,}
  \tkzTabLine{,-,z,+,d,-,z,+,}
\end{tikzpicture}
\end{center}
\explique{$S=\res{\left]-\infty\,;1\right[\cup\left]2\,;4\right[}$.}{}
\end{exo}

\begin{exo}[Tout dans un membre]
\explique{$2-x\ne 0$ : $D=\R\setminus\{2\}$. Comme on ignore le signe de $2-x$, on ne
multiplie pas : on passe tout dans un membre et on étudie le signe d'un
quotient.}{méthode 10, exemple c)}
\begin{calculs}
\dfrac{x+2}{2-x}\ge 3x+1 &\iff \dfrac{x+2}{2-x}-(3x+1)\ge 0 &&\justif{on compare à $0$}\\
&\iff \dfrac{x+2-(3x+1)(2-x)}{2-x}\ge 0 &&\justif{même dénominateur $2-x$}\\
&\iff \dfrac{x+2-\left(-3x^{2}+5x+2\right)}{2-x}\ge 0 &&\justif{$(3x+1)(2-x)=6x-3x^{2}+2-x$}\\
&\iff \dfrac{3x^{2}-4x}{2-x}\ge 0 &&\justif{parenthèse précédée de « $-$ » : on change les signes}\\
&\iff \dfrac{x(3x-4)}{2-x}\ge 0 &&\justif{facteur commun $x$}
\end{calculs}
\explique{$x$ s'annule en $0$ ; $3x-4$ en $\dfrac{4}{3}$ ; $2-x$ en $2$ (et $2-x$ est
positif \emph{avant} $2$).}{attention au facteur $2-x$ : coefficient de $x$ négatif}
\begin{center}
\begin{tikzpicture}[scale=0.8]
  \tkzTabInit[lgt=2.4,espcl=1.8]{$x$/0.8,$x$/0.8,$3x-4$/0.8,$2-x$/0.8,%
    $\frac{x(3x-4)}{2-x}$/1.1}%
    {$-\infty$,$0$,$\frac{4}{3}$,$2$,$+\infty$}
  \tkzTabLine{,-,z,+,t,+,t,+,}
  \tkzTabLine{,-,t,-,z,+,t,+,}
  \tkzTabLine{,+,t,+,t,+,z,-,}
  \tkzTabLine{,+,z,-,z,+,d,-,}
\end{tikzpicture}
\end{center}
\explique{$S=\res{\left]-\infty\,;0\right]\cup\left[\dfrac{4}{3}\,;2\right[}$.}{zéros
du numérateur compris ; $2$ exclu}
\rem{Erreur classique : multiplier par $2-x$ sans connaître son signe. Contrôle avec
$x=\frac{3}{2}$, dans $S$ : $\frac{7/2}{1/2}=7$ et $3\times\frac{3}{2}+1=\frac{11}{2}$,
et $7\ge\frac{11}{2}$. Avec $x=3$, hors de $S$ : $\frac{5}{-1}=-5$, qui n'est pas
$\ge 10$.}
\end{exo}

\partie{Problèmes}

\begin{exo}[Le nombre mystère]
\explique{Soit $x$ le nombre cherché, avec $x\ge 0$.}{\textbf{inconnue} : ce qu'elle
représente, ses valeurs possibles}
\begin{calculs}
x^{2}-3x &= 40 &&\justif{\textbf{équation} : le carré, diminué du triple}\\
x^{2}-3x-40 &= 0 &&\justif{tout dans un membre}\\
\Delta &= 9+160 &&\justif{\textbf{résolution}}\\
\Delta &= 169=13^{2} &&\\
x_{1} &= \dfrac{3+13}{2}=8 &&\\
x_{2} &= \dfrac{3-13}{2}=-5 &&
\end{calculs}
\explique{Le nombre est positif : c'est \res{$8$}.}{\textbf{retour à l'énoncé} : on
rejette $-5$ ; contrôle : $64-24=40$}
\end{exo}

\begin{exo}[Une voile d'ombrage]
\explique{Soit $x$ la longueur du plus petit côté de l'angle droit, en m, avec $x>0$ ;
l'autre mesure $x+2$.}{\textbf{inconnue}}
\begin{calculs}
\dfrac{x(x+2)}{2} &= 12 &&\justif{aire d'un triangle rectangle : produit des côtés de l'angle droit, divisé par $2$}\\
x^{2}+2x-24 &= 0 &&\justif{on multiplie par $2$ ; tout dans un membre}\\
\Delta &= 4+96=100 &&\\
x_{1} &= \dfrac{-2+10}{2}=4 &&\\
x_{2} &= \dfrac{-2-10}{2}=-6 &&
\end{calculs}
\explique{Les côtés de l'angle droit mesurent \res{$4$ m et $6$ m}.}{$x>0$ : on rejette
$-6$ ; contrôle : $\frac{4\times 6}{2}=12$}
\end{exo}

\begin{exo}[Agrandir une photo]
\explique{Soit $c$ le côté de la photo de départ, en cm, avec $c>0$. Le nouveau côté
mesure $c+12$.}{\textbf{inconnue}}
\begin{calculs}
(c+12)^{2} &= 4c^{2} &&\justif{la nouvelle aire est $4$ fois l'ancienne}\\
c^{2}+24c+144 &= 4c^{2} &&\justif{on développe}\\
-3c^{2}+24c+144 &= 0 &&\justif{tout dans un membre}\\
c^{2}-8c-48 &= 0 &&\justif{on divise par $-3$}\\
\Delta &= 64+192=256=16^{2} &&\\
c_{1} &= \dfrac{8+16}{2}=12 &&\\
c_{2} &= \dfrac{8-16}{2}=-4 &&
\end{calculs}
\explique{La photo de départ avait un côté de \res{$12$ cm}.}{$c>0$ ; contrôle :
$24^{2}=576=4\times 144$}
\rem{Plus rapide : $(c+12)^{2}=(2c)^{2}$ avec des longueurs positives, donc
$c+12=2c$.}
\end{exo}

\begin{exo}[La bordure d'une tablette]
\explique{Soit $x$ la largeur de la bordure, en cm, avec $x>0$. La bordure s'ajoute des
deux côtés : la tablette mesure $24+2x$ sur $16+2x$.}{\textbf{inconnue} et schéma}
\begin{calculs}
(24+2x)(16+2x) &= 560 &&\justif{aire de la face avant}\\
384+48x+32x+4x^{2} &= 560 &&\justif{on développe}\\
4x^{2}+80x-176 &= 0 &&\justif{tout dans un membre}\\
x^{2}+20x-44 &= 0 &&\justif{on divise par $4$}\\
\Delta &= 400+176=576=24^{2} &&\\
x_{1} &= \dfrac{-20+24}{2}=2 &&\\
x_{2} &= \dfrac{-20-24}{2}=-22 &&
\end{calculs}
\explique{La bordure mesure \res{$2$ cm} de large.}{$x>0$ ; contrôle :
$28\times 20=560$}
\rem{Erreur fréquente : écrire $24+x$ au lieu de $24+2x$.}
\end{exo}

\begin{exo}[Des vélos reconditionnés]
\cas{1)}
\begin{calculs}
B(5) &= -10\times 5^{2}+400\times 5-3000 &&\\
B(5) &= -250+2000-3000 &&\\
B(5) &= \res{-1250} &&
\end{calculs}
\begin{calculs}
B(12) &= -10\times 12^{2}+400\times 12-3000 &&\\
B(12) &= -1440+4800-3000 &&\\
B(12) &= \res{360} &&
\end{calculs}
\explique{Avec $5$ vélos, l'atelier perd $1250$ euros ; avec $12$ vélos, il gagne
$360$ euros.}{un bénéfice négatif est une perte}
\cas{2)}
\begin{calculs}
B(x) &= -10\left(x^{2}-40x+300\right) &&\justif{$-10$ en facteur : mêmes racines, calculs plus simples}\\
\Delta &= 1600-1200=400=20^{2} &&\justif{pour $x^{2}-40x+300$}\\
x_{1} &= \dfrac{40-20}{2}=10 &&\\
x_{2} &= \dfrac{40+20}{2}=30 &&
\end{calculs}
\explique{$a=-10<0$ : $B(x)>0$ entre les racines. $S=\res{\left]10\,;30\right[}$.}{et
bien dans $[0\,;40]$}
\explique{L'atelier est bénéficiaire s'il vend \res{de $11$ à $29$ vélos}.}{$x$ est un
nombre entier de vélos}
\cas{3)}
\begin{calculs}
\alpha &= -\dfrac{400}{2\times(-10)}=20 &&\justif{$a<0$ : maximum en $\alpha$}\\
\beta &= B(20)=-4000+8000-3000 &&\\
\beta &= 1000 &&
\end{calculs}
\explique{Le bénéfice est maximal pour \res{$20$ vélos} vendus ; il vaut alors
\res{$1000$ euros}.}{$20$ est le milieu des racines $10$ et $30$}
\end{exo}

\end{document}
